% Options for packages loaded elsewhere
\PassOptionsToPackage{unicode}{hyperref}
\PassOptionsToPackage{hyphens}{url}
\PassOptionsToPackage{dvipsnames,svgnames,x11names}{xcolor}
\documentclass[
  11pt,
  a4paper,
]{article}
\usepackage{xcolor}
\usepackage[margin=25mm]{geometry}
\usepackage{amsmath,amssymb}
\setcounter{secnumdepth}{-\maxdimen} % remove section numbering
\usepackage{iftex}
\ifPDFTeX
  \usepackage[T1]{fontenc}
  \usepackage[utf8]{inputenc}
  \usepackage{textcomp} % provide euro and other symbols
\else % if luatex or xetex
  \usepackage{unicode-math} % this also loads fontspec
  \defaultfontfeatures{Scale=MatchLowercase}
  \defaultfontfeatures[\rmfamily]{Ligatures=TeX,Scale=1}
\fi
\usepackage{lmodern}
\ifPDFTeX\else
  % xetex/luatex font selection
\fi
% Use upquote if available, for straight quotes in verbatim environments
\IfFileExists{upquote.sty}{\usepackage{upquote}}{}
\IfFileExists{microtype.sty}{% use microtype if available
  \usepackage[]{microtype}
  \UseMicrotypeSet[protrusion]{basicmath} % disable protrusion for tt fonts
}{}
\makeatletter
\@ifundefined{KOMAClassName}{% if non-KOMA class
  \IfFileExists{parskip.sty}{%
    \usepackage{parskip}
  }{% else
    \setlength{\parindent}{0pt}
    \setlength{\parskip}{6pt plus 2pt minus 1pt}}
}{% if KOMA class
  \KOMAoptions{parskip=half}}
\makeatother
\usepackage{longtable,booktabs,array}
\newcounter{none} % for unnumbered tables
\usepackage{calc} % for calculating minipage widths
% Correct order of tables after \paragraph or \subparagraph
\usepackage{etoolbox}
\makeatletter
\patchcmd\longtable{\par}{\if@noskipsec\mbox{}\fi\par}{}{}
\makeatother
% Allow footnotes in longtable head/foot
\IfFileExists{footnotehyper.sty}{\usepackage{footnotehyper}}{\usepackage{footnote}}
\makesavenoteenv{longtable}
\setlength{\emergencystretch}{3em} % prevent overfull lines
\providecommand{\tightlist}{%
  \setlength{\itemsep}{0pt}\setlength{\parskip}{0pt}}
\usepackage{microtype}
\usepackage{needspace}
\usepackage{fvextra}
\usepackage{adjustbox}
\AtBeginEnvironment{longtable}{\small\setlength{\tabcolsep}{3pt}}
\DefineVerbatimEnvironment{Highlighting}{Verbatim}{breaklines,commandchars=\\\{\}}
\setlength{\emergencystretch}{3em}
\widowpenalty=10000
\clubpenalty=10000
\allowdisplaybreaks
\usepackage{bookmark}
\IfFileExists{xurl.sty}{\usepackage{xurl}}{} % add URL line breaks if available
\urlstyle{same}
\hypersetup{
  pdftitle={The inverse self-commutator cost at inertia (m{,}2): a dimension-free closed formula{,} block-shift optimizers and a Pieri certificate},
  pdfauthor={Lluis Eriksson (Independent researcher)},
  colorlinks=true,
  linkcolor={blue},
  filecolor={Maroon},
  citecolor={Blue},
  urlcolor={blue},
  pdfcreator={LaTeX via pandoc}}

\title{The inverse self-commutator cost at inertia (m,2): a
dimension-free closed formula, block-shift optimizers and a Pieri
certificate}
\author{Lluis Eriksson (Independent researcher)}
\date{9 September 2026 - Internal revision v004}

\begin{document}
\maketitle

\subsection{Abstract}\label{abstract}

For a traceless Hermitian \(F\in M_d(\mathbb C)\) let
\(\kappa_d(F)=\tfrac12\min\{\|C\|_{HS}^2: CC^*-C^*C=2F\}\). We give the
exact value of \(\kappa_d\) on the whole inertia-\((m,\allowbreak{}2)\) stratum
(positive eigenvalues \(a_1\ge\dots\ge a_m>0\), two negative eigenvalues
\(-b_2\ge -b_1\), any number of zeros) for every \(m\ge2\) and every
\(d\): \(\kappa_d=\max(F_0,\allowbreak{}\dots,\allowbreak{}F_{m-1},\allowbreak{}G_1,\allowbreak{}G_3,\allowbreak{}\dots)\) for explicit
linear forms. The lower bound sums Lidskii--Wielandt and Weyl
inequalities with one Pieri-type Horn triple (Littlewood--Richardson
coefficient \(1\)); the upper bound is an explicit \(2\times2\)-block
weighted shift whose feasibility is a one-parameter interval recursion.
Both bounds are dimension-free, so the value is invariant under zero
padding, and the chamber structure is explicit. Every optimizer has rank
\(m\) or \(m+1\), the minimal optimal rank equals \(\max(n_+,\allowbreak{}n_-)=m\) in
every dimension, and rank \(m+1\) occurs exactly in one open chamber
after padding, where the optimal spectrum is not unique. By
\(F\mapsto-F\) the same formula covers inertia \((2,\allowbreak{}n)\).

\subsection{1. Introduction: antecedents and
contribution}\label{introduction-antecedents-and-contribution}

The quantity \(\kappa_d(F)\) measures the cheapest way, in
Hilbert--Schmidt norm, to realise a prescribed traceless Hermitian \(F\)
as a self-commutator \(\tfrac12(CC^*-C^*C)\); equivalently (Lemma 2.1 of
{[}FourLevel{]}) as \(\inf\|A\|_{HS}\|B\|_{HS}\) over Hermitian \(A,\allowbreak{}B\)
with \(-i[A,\allowbreak{}B]=F\). Its theory has been developed in a series of records
of the AI Research Record archive, cited here by descriptive keys (their
ARR identifiers and titles are in the reference list); everything used
from them is restated in this paper, so that it can be read without
them. {[}FourLevel{]} reduced \(\kappa_d\) to a finite linear program
over Horn inequalities (the \emph{Horn program}, Theorem 2.2 there,
restated as (2.1) below) and gave the exact four-level formula (Theorem
3.1 there)
\(\kappa_4=\max\{\lambda_1-\lambda_3,\allowbreak{}\lambda_2-\lambda_4,\allowbreak{}\lambda_1-2\lambda_2-\lambda_3,\allowbreak{}\lambda_2+2\lambda_3-\lambda_4\}\).
{[}FiveLevel{]} gave the twelve-term five-level formula (Theorem 3.1
there). {[}OneSpike{]} proved the one-spike formula
\(\kappa_d=\sum_j j\,\allowbreak{}b_j\) when one sign has multiplicity one (Theorem
3.1 there), with rigidity for balanced optima, and noted that the cost
is not known in closed form on general sign patterns. {[}RankAdaptive{]}
proved the rank-adaptive bound \(P\le\kappa_d\le(\rho/2)P\),
\(\rho=\operatorname{rank}F\) (Theorem 1.1 there), and introduced the
universal LR-one triples. {[}RankOnset{]} certified
\(r_*=r_0=\max(n_+,\allowbreak{}n_-)\) for \(d\le7\) (Theorem 3.1 there), exhibited
the first target (\(d=8\), Theorem 4.1 there) whose optimizers all have
rank above the inertia bound, and showed (Remark 3.3 there) that zero
padding can lower the value: the spectrum
\((25,\allowbreak{}18,\allowbreak{}18,\allowbreak{}-10,\allowbreak{}-17,\allowbreak{}-17,\allowbreak{}-17)/7\) has \(\kappa_7=74/7\) but
\(\kappa_8=73/7\). {[}Ceiling{]} proved the sharp inertia ceiling
\(\kappa_d\le\frac{M+1}2P\), \(M=\max(n_+,\allowbreak{}n_-)\) (Theorem 1 there).

\textbf{Contribution.} We treat the first stratum on which both signs
have multiplicity larger than one: inertia \((m,\allowbreak{}2)\) (and, by
\(F\mapsto-F\), inertia \((2,\allowbreak{}n)\)), with arbitrary zero padding. Theorem
A gives the closed formula, Theorem B its chamber structure, and Theorem
C\('\) the rank and uniqueness properties of optimizers. The proof has
two elementary halves. The lower bound is a sum of classical
Lidskii--Wielandt and Weyl inequalities plus one non-Lidskii Horn triple
whose Littlewood--Richardson coefficient is \(1\) by the Pieri rule;
only the necessity direction of the Horn--Klyachko theorem is used. The
upper bound is an explicit operator, a ``\(2\times2\)-block weighted
shift'' generalising the cumulative-sum weighted shift of
{[}RankAdaptive,OneSpike{]}, whose feasibility reduces to a
one-parameter recursion in \(2\times2\) matrices. Both halves are
independent of \(d\), which is why the value is padding-invariant here
although, as {[}RankOnset{]} showed, it is not in general; the \emph{set
of optimizers} is not padding-invariant even on this stratum (Section
7).

\textbf{Novelty, stated honestly.} Two web searches (Section 11) and the
nine records of the line found no prior closed formula for this stratum
and no prior use of block weighted shifts or of a Pieri certificate in
this problem. The nearest literature on norms in commutator
representations is the quantitative commutator theorem of
Johnson--Ozawa--Schechtman {[}JOS13{]} and its Hilbert--Schmidt version
by Angel--Schechtman {[}AS15{]}. For general factorisations \(A=[B,\allowbreak{}C]\)
of a zero-trace matrix, {[}JOS13{]} bounds the product of the two
operator norms, whereas {[}AS15{]} uses the operator norm of \(B\) and
the Hilbert--Schmidt norm of \(C\). Both objectives differ from the
self-commutator cost with prescribed spectrum studied here. These
searches do not certify priority.

\subsection{2. Setting}\label{setting}

Throughout, \(\|X\|^2=\operatorname{Tr}X^*X\) and
\(F\in M_d(\mathbb C)\) is Hermitian with \(\operatorname{Tr}F=0\) and
ordered spectrum \(\lambda_1\ge\dots\ge\lambda_d\). On the
inertia-\((m,\allowbreak{}2)\) stratum we write
\[\begin{gathered}\lambda=(a_1\ge\dots\ge a_m>0,\;0^{z},\;-b_2,\;-b_1),\\  b_1\ge b_2>0,\\  \sum_{j}a_j=b_1+b_2=:P,\\  d=m+z+2,\end{gathered}\]
and use the notation
\[\begin{gathered}A_i:=\sum_{k\ge i}a_k\ (A_{m+1}=0),\\  E_i:=a_i+a_{i+2}+a_{i+4}+\cdots\ (E_{m+1}=E_{m+2}=0),\\  g_j:=a_j-a_{j+1}\ (j\text{ odd},\ a_{m+1}:=0),\\  g^*:=\max_{j\text{ odd}\le m}g_j.\end{gathered}\]
Note \(A_i=E_i+E_{i+1}\), \(E_i\ge E_{i+1}\), and
\(E_1-E_2=\sum_{j\text{ odd}}g_j\ge g^*\).

Given \(C\) with \(CC^*-C^*C=2F\), put \(R=CC^*/2\) and \(S=C^*C/2\).
Then \(R,\allowbreak{}S\succeq0\) have the same spectrum \(s_1\ge\dots\ge s_d\ge0\)
(the \emph{common spectrum}), \(R-S=F\), and
\(\tfrac12\|C\|^2=\operatorname{Tr}R=\sum_t s_t\). Conversely any pair
\(R,\allowbreak{}S\succeq0\) with equal spectra \(R=U\Sigma U^*\), \(S=V\Sigma V^*\)
and \(R-S=F\) comes from \(C=U\sqrt{2\Sigma}V^*\). Hence
({[}FourLevel{]} Theorem 2.2, {[}RankOnset{]} Proposition 2.2, in the
\(s=p/2\) normalisation)
\[\kappa_d(F)=\min\Big\{\sum_{t<d}s_t:\ s_1\ge\dots\ge s_{d-1}\ge s_d=0,\ (s,\,-s^{\mathrm{rev}},\,\lambda)\ \text{Horn-feasible}\Big\},\tag{2.1}\]
where Horn-feasible means that \(\alpha=(s_1,\allowbreak{}\dots,\allowbreak{}s_{d-1},\allowbreak{}0)\),
\(\beta=(0,\allowbreak{}-s_{d-1},\allowbreak{}\dots,\allowbreak{}-s_1)\), \(\gamma=\lambda\) are the spectra of
Hermitian \(R\), \(-S\), \(F=R+(-S)\), i.e.~satisfy all Horn
inequalities
\(\sum_{k\in K}\gamma_k\le\sum_{i\in I}\alpha_i+\sum_{j\in J}\beta_j\),
\((I,\allowbreak{}J,\allowbreak{}K)\in T^d_r\) {[}Ho62,Kl98,KT99,Fu00{]}. The normalisation
\(s_d=0\) is legitimate because subtracting \(s_d\mathbf 1\) from \(R\)
and \(S\) preserves positivity, equal spectra and \(R-S=F\), and lowers
the cost. The sign convention is \(R-S=F\) with \(R=CC^*/2\); replacing
\(C\) by \(C^*\) swaps \(R\) and \(S\), so \(\kappa_d(-F)=\kappa_d(F)\),
which transfers every statement below to inertia \((2,\allowbreak{}n)\). Finally
\(P=\operatorname{Tr}F_+\), \(r_0(F):=\max(n_+,\allowbreak{}n_-)\) and \(r_*(F):=\)
least rank of an HS-optimizer.

In \(s\)-notation a Horn triple \((I,\allowbreak{}J,\allowbreak{}K)\) reads
\[\sum_{i\in I,\ i<d}s_i-\sum_{j\in J,\ j>1}s_{d+1-j}\ \ge\ \sum_{k\in K}\lambda_k.\tag{2.2}\]
We write \(\lambda(I)=(i_r-r,\allowbreak{}\dots,\allowbreak{}i_1-1)\) for the partition attached
to \(I=(i_1<\dots<i_r)\); by Klyachko's theorem {[}Kl98{]} the
inequality (2.2) holds for all Hermitian triples whenever the
Littlewood--Richardson coefficient
\(c^{\lambda(K)}_{\lambda(I)\lambda(J)}\) is nonzero. Only this
necessity direction is load-bearing for Theorems A and B; Horn
sufficiency (Klyachko, Knutson--Tao) enters only to identify the set of
optimal common spectra with the optimal face of the linear program
(2.1).

\subsection{3. Statements}\label{statements}

\textbf{Theorem A (inertia \((m,\allowbreak{}2)\), all \(d\)).} For every \(m\ge2\),
\(z\ge0\) and every \((a,\allowbreak{}b)\) as in Section 2,
\[\kappa_d(F)=\Phi(a;b):=\max\Big(\max_{0\le k\le m-1}F_k,\ \max_{j\ \mathrm{odd},\ 1\le j\le m}G_j\Big),\]
with
\[\begin{gathered}F_k=k\,b_1+\sum_{j=1}^m c_k(j)\,a_j,\quad c_k(j)=\begin{cases}j-k,& j\le k,\\ \lceil (j-k)/2\rceil,& j\ge k,\end{cases}\\  G_j=b_2+\sum_{i=1}^m\Big\lfloor \frac i2\Big\rfloor a_i+(a_j-a_{j+1}).\end{gathered}\]
Equivalently, in tail-sum form,
\[\begin{gathered}F_0=\sum_{t\ \mathrm{odd}}A_t,\\  F_k=b_1+\sum_{j=2}^{k}(A_j-b_2)+\sum_{t\ge k+1,\ t-k\ \mathrm{odd}}A_t\ (k\ge1),\\  G_j=b_2+\sum_{i\ \mathrm{even}<j}A_i+a_j+A_{j+2}+\sum_{i\ \mathrm{even}\ge j+3}A_i .\end{gathered}\]
The value does not depend on \(z\). The same formula gives \(\kappa_d\)
on the inertia-\((2,\allowbreak{}n)\) stratum (two positive eigenvalues
\(b_1\ge b_2\), negative eigenvalues \(-a_1,\allowbreak{}\dots,\allowbreak{}-a_n\)).

\textbf{Theorem B (chambers and exposed forms).} Let
\(k^*:=\max\{k\ge1: E_{k+1}\ge b_2\}\) (defined when \(b_2\le E_2\)) and
\(j^*\in\arg\max_{j\text{ odd}}g_j\). Then
\[\kappa_d=\begin{cases}F_{k^*}, & 0<b_2\le E_2\ \ (\text{equivalently } E_{k^*+2}\le b_2\le E_{k^*+1}),\\ F_0, & E_2\le b_2\le E_1-g^*,\\ G_{j^*}, & E_1-g^*\le b_2\ (\le b_1).\end{cases}\]
Along the \(F\)-family \(F_k\ge F_{k-1}\iff E_{k+1}\ge b_2\) (so
\(k\mapsto F_k\) is unimodal); along the \(G\)-family
\(G_j\ge G_{j'}\iff g_j\ge g_{j'}\); and
\(G_j\ge F_0\iff b_2\ge E_1-g_j\). The forms that are the unique
maximiser at some point of the stratum (``exposed'' forms) are exactly
\(F_1,\allowbreak{}\dots,\allowbreak{}F_{m-1}\), all \(G_j\) (\(j\) odd), and \(F_0\) when
\(m\ge3\); their number is \((m-1)+\lceil m/2\rceil+[m\ge3]\),
i.e.~\(2,\allowbreak{}5,\allowbreak{}6,\allowbreak{}8,\allowbreak{}9,\allowbreak{}11,\allowbreak{}12,\allowbreak{}14\) for \(m=2,\allowbreak{}\dots,\allowbreak{}9\).

\textbf{Theorem C\('\) (rank and uniqueness of optimizers).} Let \(C\)
be an HS-optimizer, i.e.~\(CC^*-C^*C=2F\) and
\(\tfrac12\|C\|^2=\kappa_d(F)\), with common spectrum \(s\).

\begin{enumerate}
\def\labelenumi{\arabic{enumi}.}
\tightlist
\item
  \(s_d=0\), \(s_{d-1}=0\) and \(s_t=0\) for \(m+2\le t\le d-1\). Hence
  \(\operatorname{rank}C\in\{m,\allowbreak{}m+1\}\), and \(r_*(F)=r_0(F)=m\) on the
  stratum in every dimension \(d\).
\item
  If \(z=0\), or the maximum \(\Phi\) is attained by some form other
  than \(G_m\) with \(m\) odd (in particular if \(m\) is even or
  \(b_2\le E_1-a_m\)), then also \(s_{m+1}=0\) and
  \(\operatorname{rank}C=m\).
\item
  If \(m\) is odd, \(z\ge1\) and \(b_2>E_1-a_m\) (the open \(G_m\)
  chamber; inside the stratum this forces \(a_m=g_m>g_j\) for every odd
  \(j<m\), so \(G_m\) is the unique maximiser), then optimizers of rank
  \(m\) and of rank \(m+1\) both exist.
\item
  The optimal common spectrum is unique for \(m=2\) (all \(z\)) and for
  \(m=3\) when \(z=0\) or when \(b_2\le E_1-a_3\); it is not unique for
  \(m=3\), \(z\ge1\), \(b_2>E_1-a_3\). For every \(m\ge4\) it is not
  unique on the open \(F_{m-3}\) chamber \(E_{m-1}<b_2<E_{m-2}\)
  whenever \(a_{m-2}>a_{m-1}\).
\end{enumerate}

In particular the value \(\kappa_d\) is padding-invariant on the stratum
but the optimizer set is not: rank-\((m+1)\) optimizers appear only
after padding.

\subsection{4. Lower bound}\label{lower-bound}

\subsubsection{4.1 Dimension-free
inequalities}\label{dimension-free-inequalities}

All inequalities below are instances of (2.2) valid in every
\(d\ge m+2\).

\textbf{Lemma 4.1 (Lidskii--Wielandt instances).} For \(I=K\) of size
\(r\) and \(J=(1,\allowbreak{}\dots,\allowbreak{}r)\), \(\lambda(J)=0\) and \(c=1\); (2.2) becomes
\(\sum_{k\in K,\allowbreak{}k<d}s_k-\sum_{t=d-r+1}^{d-1}s_t\ge\sum_{k\in K}\lambda_k\)
{[}Li50,Wi55{]}. The instances used are:

\begin{itemize}
\tightlist
\item
  \(K=(1..d-1)\): \(s_1\ge b_1\).
\item
  \((T_j)\) \(K=(j..d-1)\), \(2\le j\le m\): \(s_j\ge A_j-b_2\).
\item
  \((\Pi_t)\) \(K=(t..d-2)\), \(1\le t\le m-1\):
  \(s_t+s_{t+1}-s_{d-1}\ge A_t\).
\item
  \((B_j)\) \(K=\{j\}\cup(j+2..d-2)\), \(j+2\le m\):
  \(s_j+s_{j+2}-s_{d-1}\ge a_j+A_{j+2}\).
\item
  (Weyl) \(I=K=\{j\}\), \(J=\{1\}\), i.e.~the triple \((j;\allowbreak{}1;\allowbreak{}j)\):
  \(s_j\ge a_j\) for \(1\le j\le m\) (no other \(s\)-index occurs, since
  \(J\setminus\{1\}=\emptyset\)).
\end{itemize}

\emph{Proof.} For \(K=(t..d-2)\) the left side of (2.2) is
\((s_t+\dots+s_{d-2})-(s_{t+2}+\dots+s_{d-1})=s_t+s_{t+1}-s_{d-1}\) and
the right side is \(\sum_{k=t}^{d-2}\lambda_k=A_t\) because the zeros
contribute nothing; the other cases are identical computations.
\(\square\)

\textbf{Lemma 4.2 (lower Weyl).} \(I=(1..d-1)\),
\(J=K=[d]\setminus\{d-1\}\): \(s_2\ge b_2\).

\emph{Proof.} \(\lambda(I)=0\), \(\lambda(J)=\lambda(K)=(1)\), \(c=1\);
the sum condition \(\sum I+\sum J=\sum K+r(r+1)/2\) holds with
\(r=d-1\). The left side of (2.2) is
\((s_1+\dots+s_{d-1})-(s_1+s_3+\dots+s_{d-1})=s_2\) and the right side
is \(-\lambda_{d-1}=b_2\). It is the Weyl inequality
\(\lambda_{d-1}(F)\ge\lambda_d(R)+\lambda_{d-1}(-S)\). \(\square\)

\textbf{Lemma 4.3 (Pieri triple).} For even \(i\) with \(i+2\le d-1\)
let \(r=d-i-1\) and
\[I=\{i-1\}\cup(i+1..d-2),\qquad J=(1..d-i-2)\cup\{d-i\},\qquad K=(i..d-2).\]
Then \(\lambda(I)=((i-1)^{r-1},\allowbreak{}i-2)\), \(\lambda(J)=(1)\),
\(\lambda(K)=((i-1)^r)\), so \(c^{\lambda(K)}_{\lambda(I)\lambda(J)}=1\)
by the Pieri rule, and (2.2) reads
\[(\mathrm{Pi})_i:\qquad s_{i-1}+s_{i+2}-s_{d-1}\ \ge\ A_i ,\] which for
\(i+2=d-1\) collapses to \(s_{i-1}\ge A_i\).

\emph{Proof.} The cardinalities are \(d-i-1\) each; with \(n=d-i\),
\(\sum I+\sum J-\sum K=-1+n+\binom{n-1}{2}=\binom n2=r(r+1)/2\). The
partitions are read off from the definition (\(i_r=d-2\) gives
\(i_r-r=i-1\), and so on down to \(i_1-1=i-2\); in \(J\) only the last
index exceeds its position, by one). \(\lambda(K)/\lambda(I)\) is a
single box, so the Pieri rule gives \(c=1\) {[}Fu97{]}. The left side of
(2.2) is
\(s_{i-1}+(s_{i+1}+\dots+s_{d-2})-(s_{d-1}+\dots+s_{i+3})-s_{i+1}=s_{i-1}+s_{i+2}-s_{d-1}\),
the right side \(\sum_{k=i}^{d-2}\lambda_k=A_i\). \(\square\)

\emph{Example (\(d=6\), \(i=2\)).} Here \(r=3\), \(I=\{1,\allowbreak{}3,\allowbreak{}4\}\),
\(J=\{1,\allowbreak{}2,\allowbreak{}4\}\), \(K=\{2,\allowbreak{}3,\allowbreak{}4\}\). The associated partitions, with
trailing zeros retained, are \((1,\allowbreak{}1,\allowbreak{}0)\), \((1,\allowbreak{}0,\allowbreak{}0)\) and \((1,\allowbreak{}1,\allowbreak{}1)\).
Adding one box in the third row is the unique Pieri step, so the
coefficient is one. The sum condition is \(8+7-9=6=3\cdot4/2\). With
\(s_6=0\), the left side of (2.2) is
\[ (s_1+s_3+s_4)+(-s_6-s_5-s_3)=s_1+s_4-s_5.\] For inertia \((4,\allowbreak{}2)\)
without padding, its right side is \(a_2+a_3+a_4=A_2\). This gives the
representative inequality \(s_1+s_4-s_5\ge A_2\) with the complete
Horn/Pieri data.

The triple \((\mathrm{Pi})_i\) is not of Lidskii--Wielandt type (neither
\(I\) nor \(J\) is an initial segment \((1..r)\)). Its membership in
Fulton's recursively defined \(T^d_r\) was checked directly for
\(d\le 8\) (author's generator) and \(d\le 9\) (reviewer's independent
generator), and \(c=1\) for \(d\le12\) by both LR implementations.
Several Horn triples share the \(s\)-form \(s_{i-1}+s_{i+2}-s_{d-1}\)
--- for instance the Lidskii--Wielandt triple \(I=K=\{i-1,\allowbreak{}i+2\}\),
\(J=(1,\allowbreak{}2)\), whose right side is only \(\lambda_{i-1}+\lambda_{i+2}\);
the exhaustive enumeration of \(T^d_r\) in
{\small\texttt{pieri\_\allowbreak{}triples\_\allowbreak{}search.\allowbreak{}py}} lists all of them for \(d\le9\)
(\(6\); \(12,\allowbreak{}40\); \(32,\allowbreak{}24\); \(94,\allowbreak{}48,\allowbreak{}160\) triples for \(d=6,\allowbreak{}\dots,\allowbreak{}9\)
and \(i=2,\allowbreak{}4,\allowbreak{}\dots\)) and confirms that \((\mathrm{Pi})_i\) is the only
one whose right side is \(A_i\). That a triple outside the
Lidskii--Wielandt family is genuinely needed is visible in the linear
program: restricted to the Lidskii--Wielandt-type triples (\(I=K\),
\(J=(1..r)\) or \(J=K\), \(I=(1..r)\), which include all Weyl
inequalities), the Horn LP (2.1) falls short of \(\Phi\) by up to
\(11\%\) at \(G\)-chamber points in \(d=6,\allowbreak{}7\) (same script,
floating-point HiGHS).

\subsubsection{4.2 Certificates}\label{certificates}

A \emph{certificate} for a form is a list of the inequalities above
whose sum has every \(s_t\)-coefficient at most \(1\) on the left and
exactly the form on the right. Since \(\sum_{t<d}s_t\) minus the
certificate's left side is a nonnegative combination of the \(s_t\),
every feasible \(s\) then has cost \(\ge\) the form.

\textbf{Proposition 4.4 (certificate for \(F_k\)).}

\begin{itemize}
\tightlist
\item
  \(k=0\): add \(\Pi_1,\allowbreak{}\Pi_3,\allowbreak{}\Pi_5,\allowbreak{}\dots\), ending with \(s_m\ge a_m\)
  if \(m\) is odd. The sum is
  \(\sum_{t\le m}s_t-Ns_{d-1}\ge A_1+A_3+\dots=F_0\), \(N\) the number
  of \(\Pi\)'s.
\item
  \(k\ge1\): add \([s_1\ge b_1]\), \(T_2,\allowbreak{}\dots,\allowbreak{}T_k\), and
  \(\Pi_{k+1},\allowbreak{}\Pi_{k+3},\allowbreak{}\dots\), ending with \(s_m\ge a_m\) if \(m-k\)
  is odd. The sum is
  \(\sum_{t\le m}s_t-Ns_{d-1}\ge b_1+\sum_{j=2}^k(A_j-b_2)+\sum_{t\ge k+1,\allowbreak{}\ t-k\text{ odd}}A_t=F_k\).
\end{itemize}

Each index \(1,\allowbreak{}\dots,\allowbreak{}m\) occurs exactly once, no index above \(m\)
occurs with positive coefficient, and
\(\sum_{t<d}s_t=[\text{left side}]+\sum_{m<t<d}s_t+Ns_{d-1}\ge F_k\).

\textbf{Proposition 4.5 (certificate for \(G_j\), \(j\) odd).} Add
\([s_2\ge b_2]\), \((\mathrm{Pi})_i\) for every even \(i<j\), then
\(B_j\) (replaced by \(s_j\ge a_j\) if \(j+2>m\)), then
\(\Pi_{j+3},\allowbreak{}\Pi_{j+5},\allowbreak{}\dots\) (the last replaced by \(s_m\ge a_m\) if
needed). The indices used are \(2\);
\(\{1,\allowbreak{}4\},\allowbreak{}\{3,\allowbreak{}6\},\allowbreak{}\dots,\allowbreak{}\{j-2,\allowbreak{}j+1\}\); \(\{j,\allowbreak{}j+2\}\);
\(\{j+3,\allowbreak{}j+4\},\allowbreak{}\dots\) --- each of \(1,\allowbreak{}\dots,\allowbreak{}m\) at most once --- and the
sum is
\[\sum_{t\le m}s_t\ (+\,s_{m+1})\ -Ns_{d-1}\ \ge\ b_2+\sum_{i\text{ even}<j}A_i+a_j+A_{j+2}+\sum_{i\text{ even}\ge j+3}A_i=G_j .\]
The term \(+s_{m+1}\) is present exactly when \(j=m\) is odd and
\(z\ge1\): it comes from \((\mathrm{Pi})_{m-1}\), whose
\(s_{i+2}=s_{m+1}\) collapses against \(-s_{d-1}\) only when \(z=0\). It
is harmless for the lower bound (the coefficient is still \(\le1\)), but
it matters for rigidity (Section 7).

\emph{Proof of both propositions.} Direct addition; the index
bookkeeping is as displayed and the right sides are the tail-sum forms
of Theorem A. The tail-sum and closed forms coincide: expanding
\(A_j-b_2=b_1-(a_1+\dots+a_{j-1})\) in \(F_k\) gives coefficient
\(-(k-i)=i-k\) for \(a_i\), \(i\le k\), and
\(\#\{t\in[k+1,\allowbreak{}i]:t-k\text{ odd}\}=\lceil (i-k)/2\rceil\) for
\(i\ge k\); the count for \(G_j\) is similar (verified symbolically for
\(m\le9\) in {\small\texttt{symbolic.\allowbreak{}py}}). \(\square\)

\textbf{Corollary 4.6.} \(\kappa_d(F)\ge\Phi(a;\allowbreak{}b)\) for all \(m\ge2\),
\(z\ge0\).

\emph{Proof.} Let \(C\) be feasible with common spectrum \(s\) (not
necessarily \(s_d=0\)). The shifted pair \(R-s_d\mathbf 1\),
\(-(S-s_d\mathbf 1)\) has spectra \(\alpha,\allowbreak{}\beta\) built from
\(s'=s-s_d\) with \(s'_d=0\) and sums to \(F\), so all inequalities of
4.1 hold for \(s'\), and
\(\tfrac12\|C\|^2=\sum_{t\le d}s_t\ge\sum_{t<d}s'_t\ge\) every form.
\(\square\)

\subsubsection{\texorpdfstring{4.3 Write-outs for
\(m=4\)}{4.3 Write-outs for m=4}}\label{write-outs-for-m4}

The reviewer's independent re-derivation ({\small\texttt{my\_\allowbreak{}cert.\allowbreak{}py}}, own
Fulton generator and own LR coefficients by the bialternant formula)
produced all \(168\) certificates for \(m\le7\), \(z\in\{0,\allowbreak{}1,\allowbreak{}2,\allowbreak{}5\}\);
the case \(m=4\) at \(z=0,\allowbreak{}1,\allowbreak{}5\) (\(d=6,\allowbreak{}7,\allowbreak{}11\)) reads as follows, \(b_1\)
eliminated via \(b_1=P-b_2\).

\(F_2\):

{\def\LTcaptype{none} % do not increment counter
\begin{longtable}[]{@{}
  >{\raggedright\arraybackslash}p{(\linewidth - 6\tabcolsep) * \real{0.1176}}
  >{\raggedright\arraybackslash}p{(\linewidth - 6\tabcolsep) * \real{0.4118}}
  >{\raggedright\arraybackslash}p{(\linewidth - 6\tabcolsep) * \real{0.2059}}
  >{\raggedright\arraybackslash}p{(\linewidth - 6\tabcolsep) * \real{0.2647}}@{}}
\toprule\noalign{}
\begin{minipage}[b]{\linewidth}\raggedright
inequality
\end{minipage} & \begin{minipage}[b]{\linewidth}\raggedright
\((I;\allowbreak{}J;\allowbreak{}K)\)
\end{minipage} & \begin{minipage}[b]{\linewidth}\raggedright
\(s\)-form
\end{minipage} & \begin{minipage}[b]{\linewidth}\raggedright
right side
\end{minipage} \\
\midrule\noalign{}
\endhead
\bottomrule\noalign{}
\endlastfoot
\(s_1\ge b_1\) & \((1..d{-}1;\allowbreak{}\,\allowbreak{}1..d{-}1;\allowbreak{}\,\allowbreak{}1..d{-}1)\) & \(s_1\) &
\(a_1+a_2+a_3+a_4-b_2\) \\
\(T_2\) & \((2..d{-}1;\allowbreak{}\,\allowbreak{}1..d{-}2;\allowbreak{}\,\allowbreak{}2..d{-}1)\) & \(s_2\) &
\(a_2+a_3+a_4-b_2\) \\
\(\Pi_3\) & \((3..d{-}2;\allowbreak{}\,\allowbreak{}1..d{-}4;\allowbreak{}\,\allowbreak{}3..d{-}2)\) & \(s_3+s_4-s_{d-1}\) &
\(a_3+a_4\) \\
sum & & \(s_1+s_2+s_3+s_4-s_{d-1}\) & \(a_1+2a_2+3a_3+3a_4-2b_2=F_2\) \\
\end{longtable}
}

For \(d=6,\allowbreak{}7,\allowbreak{}11\) the three \(\Pi_3\) triples are
\(((3,\allowbreak{}4);\allowbreak{}(1,\allowbreak{}2);\allowbreak{}(3,\allowbreak{}4))\), \(((3,\allowbreak{}4,\allowbreak{}5);\allowbreak{}(1,\allowbreak{}2,\allowbreak{}3);\allowbreak{}(3,\allowbreak{}4,\allowbreak{}5))\) and
\(((3..9);\allowbreak{}(1..7);\allowbreak{}(3..9))\).

\(G_3\):

{\def\LTcaptype{none} % do not increment counter
\begin{longtable}[]{@{}
  >{\raggedright\arraybackslash}p{(\linewidth - 6\tabcolsep) * \real{0.1176}}
  >{\raggedright\arraybackslash}p{(\linewidth - 6\tabcolsep) * \real{0.4118}}
  >{\raggedright\arraybackslash}p{(\linewidth - 6\tabcolsep) * \real{0.2059}}
  >{\raggedright\arraybackslash}p{(\linewidth - 6\tabcolsep) * \real{0.2647}}@{}}
\toprule\noalign{}
\begin{minipage}[b]{\linewidth}\raggedright
inequality
\end{minipage} & \begin{minipage}[b]{\linewidth}\raggedright
\((I;\allowbreak{}J;\allowbreak{}K)\)
\end{minipage} & \begin{minipage}[b]{\linewidth}\raggedright
\(s\)-form
\end{minipage} & \begin{minipage}[b]{\linewidth}\raggedright
right side
\end{minipage} \\
\midrule\noalign{}
\endhead
\bottomrule\noalign{}
\endlastfoot
\(s_2\ge b_2\) & \((1..d{-}1;\allowbreak{}\ 1..d{-}2,\allowbreak{}d;\allowbreak{}\ 1..d{-}2,\allowbreak{}d)\) & \(s_2\) &
\(a_1+a_2+a_3+a_4-b_1\) \\
\((\mathrm{Pi})_2\) & \((1,\allowbreak{}3..d{-}2;\allowbreak{}\ 1..d{-}4,\allowbreak{}d{-}2;\allowbreak{}\ 2..d{-}2)\) &
\(s_1+s_4-s_{d-1}\) & \(a_2+a_3+a_4\) \\
Weyl & \((3;\allowbreak{}1;\allowbreak{}3)\) & \(s_3\) & \(a_3\) \\
sum & & \(s_1+s_2+s_3+s_4-s_{d-1}\) & \(b_2+a_2+2a_3+a_4=G_3\) \\
\end{longtable}
}

For \(d=6,\allowbreak{}7,\allowbreak{}11\) the Pieri triples are \(((1,\allowbreak{}3,\allowbreak{}4);\allowbreak{}(1,\allowbreak{}2,\allowbreak{}4);\allowbreak{}(2,\allowbreak{}3,\allowbreak{}4))\),
\(((1,\allowbreak{}3,\allowbreak{}4,\allowbreak{}5);\allowbreak{}(1,\allowbreak{}2,\allowbreak{}3,\allowbreak{}5);\allowbreak{}(2,\allowbreak{}3,\allowbreak{}4,\allowbreak{}5))\) and
\(((1,\allowbreak{}3,\allowbreak{}\dots,\allowbreak{}9);\allowbreak{}(1,\allowbreak{}\dots,\allowbreak{}7,\allowbreak{}9);\allowbreak{}(2,\allowbreak{}\dots,\allowbreak{}9))\). In all cases the only
index above \(m=4\) is \(d-1\) with coefficient \(-1\). The Weyl row is
the Lidskii--Wielandt instance \(I=K=\{3\}\), \(J=\{1\}\) of Lemma 4.1,
as in the reviewer's listing {\small\texttt{my\_\allowbreak{}cert.\allowbreak{}log}} (version 1 misprinted
it as \((3;\allowbreak{}3;\allowbreak{}3)\), which is not a Horn triple: \(3+3\ne3+1\)).

Each certificate is an exact dual-feasible vector of the program (2.1)
(all multipliers \(1\) in \(s\)-units, \(1/2\) in the \(p\)-units of
{[}FourLevel{]}) with dual objective equal to the form; together with
the primal points of Section 5 this shows that the LP value equals the
form at every point of its chamber without any floating-point
computation. The certificates coincide with the HiGHS dual solutions
extracted by {\small\texttt{horn\_\allowbreak{}duals\_\allowbreak{}pad.\allowbreak{}py}}.

\subsection{\texorpdfstring{5. Upper bound: \(2\times2\)-block weighted
shifts}{5. Upper bound: 2\textbackslash times2-block weighted shifts}}\label{upper-bound-2times2-block-weighted-shifts}

\subsubsection{5.1 The construction}\label{the-construction}

Partition the nonzero spectrum (optionally together with some zeros)
into ordered \emph{layers} \(\Lambda_0,\allowbreak{}\Lambda_1,\allowbreak{}\dots,\allowbreak{}\Lambda_T\) of
size \(1\) or \(2\), with \(\Lambda_0\) nonpositive and \(\Lambda_T\)
nonnegative, and let \(D_t=\operatorname{diag}\Lambda_t\). Take
\(C=\sum_{t=1}^TM_t\) block-bidiagonal with
\(M_t:\Lambda_{t-1}\to\Lambda_t\). Then
\(CC^*=\bigoplus_tM_tM_t^*=:\bigoplus 2R_t\) (with \(R_t\) on
\(\Lambda_t\)) and \(C^*C=\bigoplus_tM_t^*M_t=:\bigoplus 2S_t\) (\(S_t\)
on \(\Lambda_{t-1}\)), so \(CC^*-C^*C=2F\) iff
\[\begin{gathered}S_1=-D_0,\\  R_t-S_{t+1}=D_t\ (1\le t<T),\\  R_T=D_T,\\  R_t\ \text{and}\ S_t\ \text{have the same nonzero spectrum}.\end{gathered}\tag{5.1}\]
Conversely, given PSD \(R_t,\allowbreak{}S_t\) with equal nonzero spectra,
\(M_t=V\sqrt{2\Sigma}W^*\) realises them. The cost is
\[\tfrac12\|C\|^2=\sum_{t=1}^T\operatorname{Tr}S_t=\sum_{t=1}^T\Big(-\sum_{u<t}\operatorname{Tr}D_u\Big),\]
a sum of layer partial sums --- the block analogue of the excursion cost
of {[}RankAdaptive{]}. With singleton layers only this is the weighted
shift of {[}RankAdaptive,OneSpike{]}. Eigenvalues of \(F\) (zeros) not
placed in any layer are left outside: \(C\) vanishes on their span,
which lies in \(\ker C\cap\ker C^*\), and (5.1) on the layers is then
equivalent to \(CC^*-C^*C=2F\) on all of \(\mathbb C^d\). We use three
families of layerings:

\begin{itemize}
\tightlist
\item
  \(\mathcal L(F_0)\): \(\{-b_1,\allowbreak{}-b_2\},\allowbreak{}\{a_1,\allowbreak{}a_2\},\allowbreak{}\{a_3,\allowbreak{}a_4\},\allowbreak{}\dots\);
\item
  \(\mathcal L(F_k)\), \(k\ge1\):
  \(\{-b_1\},\allowbreak{}\{a_1\},\allowbreak{}\dots,\allowbreak{}\{a_{k-1}\},\allowbreak{}\{a_k,\allowbreak{}-b_2\},\allowbreak{}\{a_{k+1},\allowbreak{}a_{k+2}\},\allowbreak{}\dots\);
\item
  \(\mathcal L(G_j)\), \(j\) odd: \(\{-b_2\}\) followed by the
  consecutive pairs of \((a_0,\allowbreak{}a_1,\allowbreak{}\dots,\allowbreak{}a_m)\), \(a_0:=-b_1\), with
  \(a_j\) and \(a_{j+1}\) transposed:
  \(\{-b_1,\allowbreak{}a_1\},\allowbreak{}\{a_2,\allowbreak{}a_3\},\allowbreak{}\dots,\allowbreak{}\{a_{j-3},\allowbreak{}a_{j-2}\},\allowbreak{}\{a_{j-1},\allowbreak{}a_{j+1}\},\allowbreak{}\{a_j,\allowbreak{}a_{j+2}\},\allowbreak{}\{a_{j+3},\allowbreak{}a_{j+4}\},\allowbreak{}\dots\)
  (for \(j=1\): \(\{-b_1,\allowbreak{}a_2\},\allowbreak{}\{a_1,\allowbreak{}a_3\},\allowbreak{}\{a_4,\allowbreak{}a_5\},\allowbreak{}\dots\); for
  \(j=m\): the pairs of \((a_0,\allowbreak{}\dots,\allowbreak{}a_{m-2})\) followed by
  \(\{a_{m-1}\},\allowbreak{}\{a_m\}\)).
\end{itemize}

A trailing odd element forms a singleton layer. The layer partial sums
are \(A_1,\allowbreak{}A_3,\allowbreak{}\dots\) for \(\mathcal L(F_0)\);
\(b_1,\allowbreak{}\ A_2-b_2,\allowbreak{}\dots,\allowbreak{}A_k-b_2,\allowbreak{}\ A_{k+1},\allowbreak{}A_{k+3},\allowbreak{}\dots\) for
\(\mathcal L(F_k)\); and
\(b_2,\allowbreak{}\ A_2,\allowbreak{}A_4,\allowbreak{}\dots,\allowbreak{}A_{j-1},\allowbreak{}\ a_j+A_{j+2},\allowbreak{}\ A_{j+3},\allowbreak{}A_{j+5},\allowbreak{}\dots\)
for \(\mathcal L(G_j)\). Hence \(\mathrm{cost}(\mathcal L(X))=X\) for
each form \(X\) (tail-sum form; verified symbolically for \(m\le9\)). In
all three families
\(\operatorname{rank}C=\sum_t\operatorname{rank}R_t\le m\).

\subsubsection{5.2 Feasibility as an interval
recursion}\label{feasibility-as-an-interval-recursion}

\textbf{Lemma 5.0 (feasibility depends only on the spectra of the
\(S_t\)).} Fix a layering. By (5.1) the traces are forced,
\(\operatorname{Tr}S_t=-\sum_{u<t}\operatorname{Tr}D_u\) and
\(\operatorname{Tr}R_t=\operatorname{Tr}S_t\), and \(R_t\) may be
\emph{any} PSD matrix on \(\Lambda_t\) with the nonzero spectrum of
\(S_t\) (given \(S_t=W\Sigma W^*\) and such an \(R_t=V\Sigma V^*\),
\(M_t=V\sqrt{2\Sigma}W^*\) realises the pair); in particular
\(\operatorname{rank}S_t\le|\Lambda_t|\) is necessary. Hence the set of
matrices \(S_{t+1}=R_t-D_t\succeq0\) attainable at step \(t\) depends on
the past only through the spectrum of \(S_t\), and for a two-element
layer \(\Lambda_{t-1}\) the only free datum is the smaller eigenvalue
\(\sigma_2(S_t)\) (the larger one is \(\operatorname{Tr}S_t-\sigma_2\)).
A layering is feasible iff the recursion of Lemmas 5.1--5.2, started at
\(\Lambda_0\), has a nonempty set of attainable \(\sigma_2\) at every
layer and meets the terminal condition at \(\Lambda_T\); and because
Lemma 5.2 is exact in both directions, every value in each attainable
interval is realised by some solution of (5.1) --- this is used in
Section 7.

\textbf{Lemma 5.1 (step lemma).} Let \(R\succeq0\) be \(2\times2\) with
eigenvalues \(\sigma_1\ge\sigma_2\) and
\(D=\operatorname{diag}(\delta,\allowbreak{}\delta')\), \(\delta\ge\delta'\). With
\(\tau_1\ge\tau_2\) the eigenvalues of \(R-D\),
\(\tau_1+\tau_2=\sigma_1+\sigma_2-\delta-\delta'\) and, as the
eigenbasis of \(R\) rotates, \(\tau_2\) takes exactly the values in
\([\sigma_2-\delta,\allowbreak{}\ \min(\sigma_1-\delta,\allowbreak{}\sigma_2-\delta')]\).

\emph{Proof.} With \(R\) the rotation of
\(\operatorname{diag}(\sigma_1,\allowbreak{}\sigma_2)\) by angle \(\theta\) and
\(x=\cos^2\theta\),
\[\det(R-D)=(\sigma_2-\delta)(\sigma_1-\delta')+(\sigma_1-\sigma_2)(\delta-\delta')\,x,\]
affine and nondecreasing in \(x\in[0,\allowbreak{}1]\); the trace is fixed, so
\(\tau_2=\tfrac12(\tau-\sqrt{\tau^2-4\det})\) is continuous and monotone
in \(x\) and ranges over the interval between its values at \(x=0\)
(eigenvalues \(\sigma_2-\delta,\allowbreak{}\sigma_1-\delta'\)) and \(x=1\)
(eigenvalues \(\sigma_1-\delta,\allowbreak{}\sigma_2-\delta'\)). The endpoints are
the three Weyl bounds. \(\square\)

Track, for each layer, the \emph{state} \((\operatorname{Tr},\allowbreak{}[lo,\allowbreak{}hi])\):
the trace of \(S_t\) (forced:
\(\operatorname{Tr}S_t=-\sum_{u<t}\operatorname{Tr}D_u\)) and the set of
attainable second eigenvalues \(\sigma_2\) of \(S_t\) (a \(2\times2\)
matrix when \(|\Lambda_{t-1}|=2\); we always have
\(lo\le hi\le\operatorname{Tr}/2\)).

\textbf{Lemma 5.2 (interval recursion).} If \(\sigma_2\) ranges over
\([lo,\allowbreak{}hi]\) and the next layer is the pair \(\{\delta\ge\delta'\}\), the
set of attainable \(\tau_2\ge0\) for \(S_{t+1}=R_t-D_t\) is the interval
\[[lo',hi']=\Big[\max(0,\,lo-\delta),\ \min\big(hi-\delta',\ \operatorname{Tr}-\delta-lo,\ (\operatorname{Tr}-\delta-\delta')/2\big)\Big],\]
and the layer is feasible iff \(lo'\le hi'\). Transitions with
singletons: \(\Lambda_0=\{-b\}\) gives state \((b,\allowbreak{}[0,\allowbreak{}0])\) and
\(\Lambda_0=\{-b_1,\allowbreak{}-b_2\}\) gives \((b_1+b_2,\allowbreak{}[b_2,\allowbreak{}b_2])\); a singleton
\(\{\delta\}\) after a pair requires \(0\in[lo,\allowbreak{}hi]\) (rank-one \(S_t\))
and leads to the scalar \(\operatorname{Tr}-\delta\ge0\), whose
successor pair starts from \([0,\allowbreak{}0]\); the last layer
\(\{\delta,\allowbreak{}\delta'\}\) requires \(\delta'\in[lo,\allowbreak{}hi]\), and a last
singleton is automatic given the trace.

\emph{Proof.} For fixed \(\sigma=\sigma_2\) the fibre is
\([\max(0,\allowbreak{}\sigma-\delta),\allowbreak{}\min(\operatorname{Tr}-\delta-\sigma,\allowbreak{}\sigma-\delta')]\)
by Lemma 5.1, nonempty iff
\(\sigma\in[\delta',\allowbreak{}\operatorname{Tr}-\delta]\). Over the effective
domain
\([\max(lo,\allowbreak{}\delta'),\allowbreak{}\min(hi,\allowbreak{}\operatorname{Tr}-\delta,\allowbreak{}\operatorname{Tr}/2)]\)
the endpoints vary continuously and the fibres are nonempty, so their
union is an interval; its lower end is \(\max(0,\allowbreak{}lo-\delta)\) (as
\(\delta'\le\delta\)) and its upper end is the maximum of the concave
function
\(\sigma\mapsto\min(\operatorname{Tr}-\delta-\sigma,\allowbreak{}\sigma-\delta')\)
over the domain, which is the minimum of the three displayed quantities.
If \(lo'\le hi'\) then \(hi'\ge0\) forces \(hi\ge\delta'\),
\(lo\le\operatorname{Tr}-\delta\) and
\(\delta'\le\operatorname{Tr}-\delta\), so the effective domain is
nonempty and the formula is exact. The singleton rules are the rank-one
cases of (5.1). \(\square\)

The recursion is implemented exactly (rational arithmetic) in
{\small\texttt{layer\_\allowbreak{}chain.\allowbreak{}feasible}} and, independently with the explicit
domain intersection, in the reviewer's {\small\texttt{my\_\allowbreak{}construct.\allowbreak{}forward}};
the two agree on all \(2708\) boundary-heavy exact points of
{\small\texttt{adversarial.\allowbreak{}py}}.

\textbf{Lemma 5.3 (pair chains from \(lo=0\)).} Consider a state
\((\operatorname{Tr},\allowbreak{}[0,\allowbreak{}h])\) followed by pair layers
\(\{\delta_1\ge\delta'_1\},\allowbreak{}\dots,\allowbreak{}\{\delta_q\ge\delta'_q\}\) of
nonnegative numbers, possibly followed by singleton layers of
nonnegative numbers, with \(\operatorname{Tr}\) the total of all these
layers. The chain is feasible iff \(h\ge\sum_{i=1}^q\delta'_i\).

\emph{Proof.} Since \(lo=0\) propagates (\(\max(0,\allowbreak{}0-\delta)=0\)), write
\(h_1=h\) and
\(h_{i+1}=\min(h_i-\delta'_i,\allowbreak{}\ \operatorname{Tr}_i-\delta_i,\allowbreak{}\ (\operatorname{Tr}_i-\delta_i-\delta'_i)/2)\),
\(\operatorname{Tr}_i\) the trace before pair \(i\). Let
\(N_i=\sum_{i'\ge i}\delta'_{i'}\). If \(h_i\ge N_i\) then
\(h_{i+1}\ge N_{i+1}\): the first term by hypothesis;
\(\operatorname{Tr}_i-\delta_i\) is the sum of everything after
\(\delta_i\), which contains all later \(\delta'\); and
\((\operatorname{Tr}_i-\delta_i-\delta'_i)/2\ge N_{i+1}\) because the
later layers contain, besides the \(\delta'_{i'}\), the
\(\delta_{i'}\ge\delta'_{i'}\). Hence all \(h_i\ge0\), and the terminal
condition (\(\delta'_q\in[0,\allowbreak{}h_q]\) for a pair end, \(h_{q+1}\ge0\) for a
singleton end) holds. Conversely \(h_{i+1}\le h_i-\delta'_i\) gives
\(h_q\le h-\sum_{i<q}\delta'_i\), so the terminal condition forces
\(h\ge N_1\). \(\square\)

\subsubsection{5.3 Closed-form feasibility
regions}\label{closed-form-feasibility-regions}

\textbf{Proposition 5.4.} Inside the stratum (\(0<b_2\le b_1\)):

\begin{enumerate}
\def\labelenumi{\arabic{enumi}.}
\tightlist
\item
  \(\mathcal L(F_k)\), \(k\ge1\), is feasible iff
  \(E_{k+2}\le b_2\le E_{k+1}\).
\item
  \(\mathcal L(G_j)\) is feasible iff \(b_2\ge E_1-g_j\); inside the
  stratum this forces \(g_j=g^*\).
\item
  \(\mathcal L(F_0)\) is feasible iff \(E_2\le b_2\le E_1-g^*\).
\end{enumerate}

\emph{Proof of 1.} The singleton chain
\(\{-b_1\},\allowbreak{}\{a_1\},\allowbreak{}\dots,\allowbreak{}\{a_{k-1}\}\) has scalars
\(b_1,\allowbreak{}\ b_1-a_1,\allowbreak{}\dots,\allowbreak{}\ A_k-b_2\), nonnegative iff \(b_2\le A_k\). The
mixed layer \(\{a_k,\allowbreak{}-b_2\}\) (\(\delta=a_k\), \(\delta'=-b_2\)) starts
from \((A_k-b_2,\allowbreak{}[0,\allowbreak{}0])\) and by Lemma 5.2 yields \((A_{k+1},\allowbreak{}[0,\allowbreak{}h])\)
with \(h=\min(b_2,\allowbreak{}\ A_{k+1}-b_2,\allowbreak{}\ A_{k+1}/2)=\min(b_2,\allowbreak{}A_{k+1}-b_2)\).
The remaining pairs \(\{a_{k+1},\allowbreak{}a_{k+2}\},\allowbreak{}\dots\) have
\(\sum\delta'=E_{k+2}\), so by Lemma 5.3 the layering is feasible iff
\(\min(b_2,\allowbreak{}A_{k+1}-b_2)\ge E_{k+2}\), i.e.~iff
\(E_{k+2}\le b_2\le A_{k+1}-E_{k+2}=E_{k+1}\); this implies
\(b_2\le A_k\), so the singleton chain is nonnegative whenever the whole
layering is feasible. Boundary cases: for \(k=m-1\) there is no pair
after the mixed layer, only the singleton \(\{a_m\}\), and Lemma 5.3
with \(q=0\) gives feasibility iff \(h\ge0\), i.e.~iff
\(b_2\le A_m=E_m\), in agreement with \(E_{m+1}=0\); when \(m-k\) is odd
the chain ends with the singleton \(\{a_m\}\), covered by Lemma 5.3.
(When \(b_2=A_{k+1}/2\) the third term of the recursion, not the second,
is the binding one after the first pair; it never causes infeasibility,
which is all that is needed.)

\emph{Proof of 2.} Let \(j\ge3\). After \(\{-b_2\}\) the state is
\((b_2,\allowbreak{}[0,\allowbreak{}0])\); the pair \(\{-b_1,\allowbreak{}a_1\}\) (\(\delta=a_1\),
\(\delta'=-b_1\)) gives \((A_2,\allowbreak{}[0,\allowbreak{}h])\) with
\(h=\min(b_1,\allowbreak{}\ b_2-a_1,\allowbreak{}\ (P-a_1)/2)=b_2-a_1\), the other two terms being
dominated because \(b_2\le b_1\). The later pairs have
\(\delta'\)-values
\(a_3,\allowbreak{}a_5,\allowbreak{}\dots,\allowbreak{}a_{j-2},\allowbreak{}\ a_{j+1},\allowbreak{}\ a_{j+2},\allowbreak{}\ a_{j+4},\allowbreak{}\dots\), which sum
to \(E_3-(a_j-a_{j+1})=E_3-g_j\) (for \(j=m\) the trailing
\(\{a_{m-1}\},\allowbreak{}\{a_m\}\) are singletons and the sum is \(E_3-a_m\)). The
tail of the chain is: for \(m\) odd and \(j\le m-2\), pairs only (the
last pair is \(\{a_{m-1},\allowbreak{}a_m\}\)); for \(m\) even and \(j\le m-3\),
pairs followed by the singleton \(\{a_m\}\); for \(j=m-1\) (\(m\) even),
the pairs end with \(\{a_{m-2},\allowbreak{}a_m\}\) and the singleton \(\{a_{m-1}\}\)
follows; for \(j=m\) (\(m\) odd), the pairs end with
\(\{a_{m-3},\allowbreak{}a_{m-2}\}\) and the two singletons \(\{a_{m-1}\},\allowbreak{}\{a_m\}\)
follow (the scalar after the first singleton is \(a_m\ge0\), the last
singleton is automatic). In each case the \(\delta'\)-sum is \(E_3-g_j\)
as displayed, the layer entries are nonnegative and the trace is the
total, so Lemma 5.3 applies: feasible iff \(b_2-a_1\ge E_3-g_j\),
i.e.~\(b_2\ge E_1-g_j\). For \(j=1\) the first pair is \(\{-b_1,\allowbreak{}a_2\}\),
\(h=\min(b_1,\allowbreak{}\ b_2-a_2,\allowbreak{}\ (P-a_2)/2)=b_2-a_2\) by the same argument, the
later \(\delta'\)-sum is \(E_3\), and the condition is
\(b_2\ge a_2+E_3=E_1-g_1\) (for \(m=2\) the tail is the single singleton
\(\{a_1\}\) and \(E_3=0\)). Finally
\(E_1-g_j\le b_2\le P/2=(E_1+E_2)/2\) gives \(\sum_{j'\ne j}g_{j'}<g_j\)
or equality, so \(g_j=g^*\).

\emph{Proof of 3.} The initial state is \((A_1,\allowbreak{}[b_2,\allowbreak{}b_2])\) (this uses
\(b_2\le b_1\)). Let \((A_{2t+1},\allowbreak{}[lo_t,\allowbreak{}hi_t])\) be the state before the
pair \(\{a_{2t+1},\allowbreak{}a_{2t+2}\}\). By Lemma 5.2,
\(lo_t=\max\big(0,\allowbreak{}\ b_2-(a_1+a_3+\dots+a_{2t-1})\big)\) and
\[hi_t=\min\left\{\begin{gathered}b_2-(a_2+\dots+a_{2t}),\\ A_{2r}-lo_{r-1}-(a_{2r+2}+\dots+a_{2t}),\\ \tfrac12A_{2r+1}-(a_{2r+2}+\dots+a_{2t})\ \ (1\le r\le t)\end{gathered}\right\}.\]
\emph{Induction for the displayed endpoints.} Empty sums are zero, and
all dotted sums above advance by two. At \(t=0\) the minimum contains
only \(b_2\), so \(lo_0=hi_0=b_2\). Applying Lemma 5.2 to the next
nonterminal pair gives \[\begin{aligned}
lo_{t+1}&=\max(0,lo_t-a_{2t+1}),\\
hi_{t+1}&=\min\{hi_t-a_{2t+2},\ A_{2t+2}-lo_t,\ \tfrac12 A_{2t+3}\}.
\end{aligned}\] The first identity proves the formula for \(lo_{t+1}\).
Subtracting \(a_{2t+2}\) from every candidate in \(hi_t\) appends that
term to each even-indexed tail sum; the other two candidates are
precisely the displayed terms with \(r=t+1\) and an empty tail. This
proves the closed form by induction at every retained state.
Equivalently the formulas can be unfolded algebraically before testing
feasibility; the terminal and intermediate checks below show that none
of the required intervals is empty on the stated region.

\emph{Terminal conditions.} For \(m\) even the last layer is the pair
\(\{a_{m-1},\allowbreak{}a_m\}\) and requires \(lo_{q-1}\le a_m\le hi_{q-1}\),
\(q=m/2\); for \(m\) odd the last layer is the singleton \(\{a_m\}\) and
requires \(lo_q=0\le hi_q\), \(q=(m-1)/2\). The \(lo\) condition is
\(b_2\le a_1+a_3+\dots+a_{m-3}+a_m=E_1-g_{m-1}\) (\(m\) even) or
\(b_2\le a_1+\dots+a_{m-2}=E_1-g_m\) (\(m\) odd). In the \(hi\)
condition the first term gives \(b_2\ge E_2\); the \(r\)-th middle term
gives, after simplifying
\(A_{2r}-(a_{2r+2}+a_{2r+4}+\dots)=a_{2r}+E_{2r+1}\), the condition
\(b_2\le a_1+a_3+\dots+a_{2r-3}+a_{2r}+E_{2r+1}=E_1-g_{2r-1}\) for
\(j=2r-1=1,\allowbreak{}3,\allowbreak{}\dots\) (\(r\le q-1\) for \(m\) even, \(r\le q\) for \(m\)
odd); the third terms are automatic since
\(A_{2r+1}=E_{2r+1}+E_{2r+2}\ge2E_{2r+2}\). Together the terminal
conditions are exactly \(E_2\le b_2\le E_1-g_j\) for every odd
\(j\le m\), and they are necessary. \emph{Intermediate conditions.}
\(lo_t\le hi_t\) for \(t\) before the terminal index. If \(lo_t=0\) each
term of \(hi_t\) exceeds the corresponding terminal term, so the
terminal conditions suffice. If \(lo_t=b_2-(a_1+\dots+a_{2t-1})>0\) then
also \(lo_{r-1}>0\) for \(r\le t\) (\(lo\) is nonincreasing), and the
three conditions become: for the first term,
\(a_1+a_3+\dots+a_{2t-1}\ge a_2+\dots+a_{2t}\), automatic; for the
middle term with index \(r\), after substituting \(lo_{r-1}\),
\[A_{2r}-(a_{2r+2}+\dots+a_{2t})+(a_1+\dots+a_{2r-3})+(a_1+\dots+a_{2t-1})\ \ge\ 2b_2,\]
whose left side minus \(P\) equals
\(\sum_{i<r}(a_{2i-1}-a_{2i})+\sum_{i=r}^{t-1}(a_{2i+1}-a_{2i+2})\ge0\),
so it holds because \(2b_2\le P\); for the third term, put
\(O_t=\sum_{i=1}^t a_{2i-1}\). Since
\(A_{2r+1}=P-\sum_{i=1}^r(a_{2i-1}+a_{2i})\), the exact identity is
\[\begin{aligned}
\frac12A_{2r+1}-\sum_{i=r+1}^t a_{2i}+O_t-b_2
&=\frac P2-b_2+\frac12\sum_{i=1}^r(a_{2i-1}-a_{2i})\\
&\quad+\sum_{i=r+1}^t(a_{2i-1}-a_{2i})\ \ge0.
\end{aligned}\] The first term is nonnegative because \(P=b_1+b_2\) and
\(b_1\ge b_2\); the remaining differences are nonnegative by the
ordering of the \(a_i\). Empty sums are zero, so this includes \(r=t\)
and all equality cases. Thus the third candidate for \(hi_t\) is at
least \(lo_t=b_2-O_t\). Thus the intermediate conditions in the case
\(lo_t>0\) are consequences of the stratum condition \(b_2\le b_1\) and
not of the terminal conditions alone; on the stratum they are harmless.
(For \(m=2\) there is no intermediate step: the single pair
\(\{a_1,\allowbreak{}a_2\}\) is terminal and requires
\(b_2=lo_0\le a_2\le hi_0=b_2\), i.e.~\(b_2=a_2=E_2=E_1-g_1\), the
degenerate \(F_0\) region.) \(\square\)

All three regions were re-derived independently by the reviewer's
recursion ({\small\texttt{my\_\allowbreak{}construct.\allowbreak{}forward}}) and are checked against the
exact recursion on a fine rational grid with every tie and boundary
case, \(m\le10\) ({\small\texttt{coverage\_\allowbreak{}exact.\allowbreak{}py}}, Section 9).

\textbf{Proposition 5.5 (coverage).} For every \((a,\allowbreak{}b)\) in the stratum
some layering is feasible, and its cost is \(\Phi(a;\allowbreak{}b)\).

\emph{Proof.} Since \(E_{m+1}=0<b_2\): if \(b_2\le E_2\) then
\(k^*=\max\{k\ge1:E_{k+1}\ge b_2\}\le m-1\) exists and
\(E_{k^*+2}\le b_2\le E_{k^*+1}\), so \(\mathcal L(F_{k^*})\) is
feasible; if \(E_2\le b_2\le E_1-g^*\), \(\mathcal L(F_0)\) is feasible;
if \(b_2\ge E_1-g^*\), \(\mathcal L(G_{j^*})\) is feasible. The three
intervals cover \((0,\allowbreak{}\infty)\) because \(E_2\le E_1-g^*\). The cost of
the feasible layering is its form, which is \(\le\Phi\le\kappa_d\) by
Corollary 4.6, while the layering shows \(\kappa_d\le\) its cost.
\(\square\)

\emph{Proof of Theorem A.} Corollary 4.6 and Proposition 5.5 give
\(\kappa_d=\Phi\) in \(d=m+z+2\) for every \(z\); neither depends on
\(z\). The \((2,\allowbreak{}n)\) case follows from \(\kappa_d(-F)=\kappa_d(F)\).
\(\square\)

Because the regions in Proposition 5.4 are exact (both directions), the
feasibility region of each layering coincides with the chamber of its
form in Theorem B, and exactly the argmax layerings are feasible at
every point --- a fact checked at \(400\) exact points per \(m\le8\),
\(300\) per \(m\le13\), and at \(9567\) grid points for \(m\le10\) that
include every breakpoint \(E_k\), \(E_1-g_j\), \(b_1=b_2\), flat and
tied spectra ({\small\texttt{coverage\_\allowbreak{}exact.\allowbreak{}py}}). The end-to-end construction
of \(C\) (backward choice of the \(\sigma_2\)'s, rotation angles from
the affine determinant, SVD blocks) gives residuals \(\|CC^*-C^*C-2F\|\)
and \(|\tfrac12\|C\|^2-\Phi|\) at the \(10^{-14}\) level away from tie
points and \(\le2\cdot10^{-8}\) at exact ties, with
\(\operatorname{rank}C=m\).

\subsection{6. Chambers and exposedness: proof of Theorem
B}\label{chambers-and-exposedness-proof-of-theorem-b}

From the tail-sum forms, \(F_1-F_0=b_1-E_1=E_2-b_2\) and, for \(k\ge2\),
\(F_k-F_{k-1}=(A_k-b_2)+\sum_{t\ge k+1,\allowbreak{}\,\allowbreak{}t-k\text{ odd}}A_t-\sum_{t\ge k,\allowbreak{}\,\allowbreak{}t-k\text{ even}}A_t=-b_2+a_{k+1}+a_{k+3}+\dots=E_{k+1}-b_2\).
Since \(E_{k+1}\) is nonincreasing in \(k\), \(k\mapsto F_k\) increases
while \(E_{k+1}\ge b_2\) and decreases afterwards: unimodal with maximum
at \(k^*\) when \(b_2\le E_2\) and at \(k=0\) when \(b_2\ge E_2\). From
the closed forms \(G_j-G_{j'}=g_j-g_{j'}\), and since
\(F_0=\sum_i\lceil i/2\rceil a_i=\sum_i\lfloor i/2\rfloor a_i+E_1\),
\(G_j-F_0=b_2-E_1+g_j\). Hence for \(b_2\le E_2\) the maximum is
\(F_{k^*}\) (as \(E_2\le E_1-g^*\) gives \(F_0\ge G_j\) for all \(j\));
for \(E_2\le b_2\le E_1-g^*\) it is \(F_0\); for \(b_2\ge E_1-g^*\) it
is \(G_{j^*}\). This proves the chamber description.

\emph{Exposedness.} \(F_k\) (\(1\le k\le m-1\)) is the unique maximiser
on the open interval \(E_{k+2}<b_2<E_{k+1}\): there
\(F_{k'}-F_{k'-1}=E_{k'+1}-b_2\) is \(>0\) for \(k'\le k\) and \(<0\)
for \(k'>k\), and \(F_0>G_j\) for every \(j\) because
\(b_2<E_2\le E_1-g_j\). The interval lies inside the stratum because
\(E_{k+1}\le E_2\le P/2\), and it is nonempty iff \(E_{k+1}>E_{k+2}\).
Here the correct alternating-sum identities are
\(E_{k+1}-E_{k+3}=a_{k+1}\) and
\(E_{k+1}-E_{k+2}=(a_{k+1}-a_{k+2})+(a_{k+3}-a_{k+4})+\cdots\), the last
term being \(a_m>0\) when \(m-k\) is odd (version 1 wrote
\(E_{k+1}-E_{k+2}=a_{k+1}\), which is false: for \(a=(4,\allowbreak{}3,\allowbreak{}2,\allowbreak{}1)\) and
\(k=1\), \(E_2-E_3=2\ne a_2=3\)). So the open \(F_k\) chamber is
nonempty whenever \(a_{k+1}>a_{k+2}\) or \(m-k\) is odd, and empty
exactly when \(m-k\) is even and every pair satisfies
\(a_{k+1+2t}=a_{k+2+2t}\) for \(0\le t<(m-k)/2\); an exposing point
exists for every \(k\), e.g.~for \(a=(m,\allowbreak{}m-1,\allowbreak{}\dots,\allowbreak{}1)\), where
\(E_{k+1}-E_{k+2}\ge1\). \(F_0\) is the unique maximiser on
\(E_2<b_2<E_1-g^*\), nonempty iff \(\sum_{j\text{ odd}}g_j>g^*\),
i.e.~iff at least two odd gaps are positive, possible iff \(m\ge3\); for
\(m=2\), \(E_2=E_1-g_1=a_2\) and at that point \(F_0=F_1=G_1\), so
\(F_0\) is never a unique maximiser (indeed \(F_0\le\max(F_1,\allowbreak{}G_1)\)
since \(F_0-F_1=b_2-a_2\) and \(F_0-G_1=a_2-b_2\)). \(G_j\) is the
unique maximiser when \(b_2>E_1-g_j\) and \(g_j>g_{j'}\) for
\(j'\ne j\); a point with \(E_1-g_j<b_2\le P/2\) exists iff
\(\sum_{j'\ne j}g_{j'}<g_j\), which is realised by \(a_1=\dots=a_j=1\),
\(a_{j+1}=\dots=a_m=\varepsilon\) (all \(a_i=1\) when \(j=m\)). This
gives the count \((m-1)+\lceil m/2\rceil+[m\ge3]\); exact unique-argmax
sampling (\(20000\) rational points per \(m\), {\small\texttt{rank\_\allowbreak{}more.\allowbreak{}py}})
and the LP-dual counts in {\small\texttt{forms\_\allowbreak{}m*\_\allowbreak{}pad*.\allowbreak{}json}} return the same
numbers \(2,\allowbreak{}5,\allowbreak{}6,\allowbreak{}8,\allowbreak{}9,\allowbreak{}11,\allowbreak{}12,\allowbreak{}14\). \(\square\)

For example, \(a=(4,\allowbreak{}3,\allowbreak{}3,\allowbreak{}1,\allowbreak{}1)\) and \(k=1\) give \(E_2=E_3=4\): the open
\(F_1\) chamber is empty although the tail \((3,\allowbreak{}3,\allowbreak{}1,\allowbreak{}1)\) is not
constant. The corrected criterion follows because \(E_{k+1}-E_{k+2}\) is
a sum of nonnegative pair differences, plus the positive unpaired last
term when \(m-k\) is odd. It does not change the formula for the cost or
the global count of exposed forms, which asserts existence over varying
spectra.

\emph{When are the \(G\)-forms active?} Since
\(G_{j^*}\ge F_0\iff b_2\ge E_1-g^*\) and \(b_2\le P/2=(E_1+E_2)/2\) on
the stratum, some \(G_j\) attains \(\Phi\) at some point of the stratum
iff \(E_1-g^*\le P/2\), i.e.~iff \(\sum_{j\ne j^*}g_j\le g^*\) (the
largest odd gap is at least the sum of the others); a \(G_j\) is the
\emph{unique} maximiser somewhere iff this inequality is strict. When it
holds with equality, the \(G\)-forms are active exactly on the line
\(b_1=b_2\), where they tie with \(F_0\): for \(m=3\), \(a=(2,\allowbreak{}1,\allowbreak{}1)\),
\(b_1=b_2=2\) one has \(g_1=g_3=1\) and \(F_0=G_1=G_3=5\). In
particular, if all \(\lceil m/2\rceil\) odd gaps are equal to some
\(g\): for \(m\in\{3,\allowbreak{}4\}\) (two odd gaps) the \(G\)-forms are active
only at \(b_1=b_2\), tied with \(F_0\); for \(m\ge5\) and \(g>0\) they
are never active; and if \(g=0\) (\(m\) even, \(a_{2i-1}=a_{2i}\) for
all \(i\)) then \(E_1=E_2=P/2\) and every \(G_j\) ties with \(F_0\) at
\(b_1=b_2\). Version 1 stated that the \(G\)-forms are never active when
all odd gaps are equal and \(m\ge3\); that statement omitted these tie
cases.

\subsection{\texorpdfstring{7. Rank and uniqueness: proof of Theorem
C\('\)}{7. Rank and uniqueness: proof of Theorem C\textquotesingle{}}}\label{rank-and-uniqueness-proof-of-theorem-c}

\textbf{Inertia bound.} For any feasible \(C\) and \(x\in\ker C^*\),
\(\langle 2Fx,\allowbreak{}x\rangle=-\|Cx\|^2\le0\), so \(F\preceq0\) on a subspace
of dimension \(d-\operatorname{rank}C\), whence
\(\operatorname{rank}C\ge n_+=m\); symmetrically
\(\operatorname{rank}C\ge n_-\). This is the inertia obstruction, Lemma
2.1 of {[}RankOnset{]}.

\textbf{(1) and (2).} Let \(C\) be an optimizer with common spectrum
\(s\). Then
\(\sum_{t\le d}s_t=\sum_{t<d}(s_t-s_d)+d\,\allowbreak{}s_d\ge\Phi+d\,\allowbreak{}s_d\) by
Corollary 4.6 applied to \(s-s_d\), so \(s_d=0\). Let \(X\) be any form
attaining \(\Phi\) and consider its certificate, whose left side is
\(L(s)=\sum_{t\in U}s_t-Ns_{d-1}\) with \(U\subseteq\{1,\allowbreak{}\dots,\allowbreak{}m+1\}\)
the indices used with coefficient \(+1\) and \(N\ge0\) the number of
\(-s_{d-1}\) terms (possibly \(N=0\), e.g.~for \(F_{m-1}\)). Then
\(\sum_{t<d}s_t-L(s)=\sum_{t<d,\allowbreak{}\ t\notin U}s_t+Ns_{d-1}\ge0\) and
\(L(s)\ge X=\Phi=\sum_{t<d}s_t\), so every coordinate \(s_t\) with
\(t<d\), \(t\notin U\) vanishes (and \(s_{d-1}=0\) also when \(N\ge1\)).
For every \(F_k\) and every \(G_j\) with \(j<m\), and for \(G_m\) with
\(z=0\), the used indices are \(\subseteq\{1,\allowbreak{}\dots,\allowbreak{}m\}\) together with
\(-s_{d-1}\): hence \(s_t=0\) for \(m<t<d\) and
\(\operatorname{rank}C=m\). For \(G_m\) with \(m\) odd and \(z\ge1\) the
used indices are \(\{1,\allowbreak{}\dots,\allowbreak{}m+1\}\) and \(-s_{d-1}\) with
\(d-1\ge m+2\): hence \(s_t=0\) for \(m+2\le t\le d-1\), and \(s_{m+1}\)
is unconstrained by this argument. In all cases \(d-1\notin U\) (as
\(d-1\ge m+1\), with equality only when \(z=0\), where the \(s_{m+1}\)
term collapsed), so \(s_{d-1}=0\) and
\(\operatorname{rank}C=\#\{t:s_t>0\}\in\{m,\allowbreak{}m+1\}\). If some form other
than \(G_m\) attains \(\Phi\) its certificate forces \(s_{m+1}=0\); this
covers \(m\) even, \(z=0\), and every point where \(G_m\) is not the
unique maximiser, in particular \(b_2\le E_1-a_m\) (where
\(G_m\le F_0\)). The block construction has rank \(m\), so \(r_*=r_0=m\)
in every \(d\), extending the \(d\le7\) certificate of {[}RankOnset{]}
(Theorem 3.1 there) to the whole stratum. The two cases (2) and (3) are
exhaustive: for \(m\) odd and \(z\ge1\), either \(b_2\le E_1-a_m\),
where \(G_m\le F_0\) and (2) applies, or \(b_2>E_1-a_m\), the open
\(G_m\) chamber of (3).

\textbf{(3) Rank-\((m+1)\) optimizers.} The preceding certificate leaves
\(s_{m+1}\) free in this case; the following construction determines
whether positive values are attainable. Let \(m\) be odd, \(z\ge1\), and
use the layering
\[\mathcal L'_m:\qquad \{-b_2\},\ \{-b_1,a_1\},\ \{a_2,a_3\},\ \dots,\ \{a_{m-3},a_{m-2}\},\ \{a_{m-1},0\},\ \{a_m\},\]
i.e.~\(\mathcal L(G_m)\) with one zero eigenvalue placed inside the
penultimate layer; the remaining zeros are left outside (they span
\(\ker C\cap\ker C^*\)). Its cost is still \(G_m\). Number the layers
\(\Lambda_0=\{-b_2\}\), \(\Lambda_1=\{-b_1,\allowbreak{}a_1\}\), \(\dots\),
\(\Lambda_{T-1}=\{a_{m-1},\allowbreak{}0\}\), \(\Lambda_T=\{a_m\}\), \(T=(m+3)/2\).
By Lemma 5.3 (the \(\delta'\)-values after the first pair are
\(a_3,\allowbreak{}\dots,\allowbreak{}a_{m-2},\allowbreak{}0\)) it is feasible iff \(b_2-a_1\ge E_3-a_m\),
i.e.~\(b_2\ge E_1-a_m\), the closed \(G_m\) region. Its rank is
\(\sum_t\operatorname{rank}R_t=\sum_t\operatorname{rank}S_t\), where
\(\operatorname{rank}S_1=\operatorname{rank}S_T=1\) (\(S_1=b_2\);
\(S_T\) must have rank one because \(\Lambda_T\) is a singleton) and
\(\operatorname{rank}S_t=1+[\sigma_2(S_t)>0]\) for \(2\le t\le T-1\).
For \(2\le t\le T-2\) the layer \(\Lambda_t\) following \(S_t\) is a
pair \(\{a_{2t-2},\allowbreak{}a_{2t-1}\}\) with \(\delta'_t=a_{2t-1}>0\), and the
fibre of Lemma 5.2 is nonempty only when
\(\sigma_2(S_t)\ge\delta'_t>0\); so these \(S_t\) have rank \(2\) in
every solution, and
\[\operatorname{rank}C=1+\sum_{t=2}^{T-2}2+\operatorname{rank}S_{T-1}+1=m+[\sigma_2(S_{T-1})>0],\]
where \(S_{T-1}\) lives on \(\{a_{m-3},\allowbreak{}a_{m-2}\}\) (on \(\{-b_1,\allowbreak{}a_1\}\)
when \(m=3\)). Its successor \(\{a_{m-1},\allowbreak{}0\}\) is followed by the
singleton \(\{a_m\}\), which needs \(\tau_2(S_T)=0\); from
\(\sigma_2(S_{T-1})=\sigma\) this is attainable iff
\(\sigma-a_{m-1}\le0\le\min(a_m-\sigma,\allowbreak{}\sigma)\) (Lemma 5.1 with
\(\operatorname{Tr}S_{T-1}=a_{m-1}+a_m\)), i.e.~iff
\(\sigma\in[0,\allowbreak{}a_m]\). So rank \(m+1\) is attainable iff the upper end
\(h\) of the \(\sigma_2\)-interval of \(S_{T-1}\) is positive, and
\(\sigma=0\) always gives rank \(m\). When \(b_2>E_1-a_m\) the induction
of Lemma 5.3 runs with strict inequalities (\(h_1>N_1\), and every later
term exceeds \(N_{i+1}\) by at least \(a_m>0\) or by
\(\tfrac12(\sum(\delta-\delta')+a_m)>0\)), so \(h>N=0\). When
\(b_2=E_1-a_m\), \(h\le h_1-\sum_{i<q}\delta'_i=(b_2-a_1)-(E_3-a_m)=0\),
consistent with (2) since \(G_m=F_0\) there. This proves (3); both
directions are checked exactly on the grid of
{\small\texttt{coverage\_\allowbreak{}exact.\allowbreak{}py}} for odd \(m\le9\).

\textbf{Example (the reviewer's counterexample to the original rank
claim).} \(\lambda=(2,\allowbreak{}2,\allowbreak{}\tfrac32,\allowbreak{}0,\allowbreak{}-\tfrac52,\allowbreak{}-3)\): \(m=3\), \(z=1\),
\(a=(2,\allowbreak{}2,\allowbreak{}\tfrac32)\), \(b=(3,\allowbreak{}\tfrac52)\), \(g_1=0\), \(g_3=\tfrac32\),
\(E_1-g_3=2<b_2=\tfrac52\le b_1\). The forms are \(F_0=7\),
\(F_1=\tfrac{13}2\), \(F_2=\tfrac{11}2\), \(G_1=6\),
\(G_3=\tfrac{15}2=\kappa_6\). The \(G_3\) certificate is
\([s_2\ge b_2]+(\mathrm{Pi})_2+[s_3\ge a_3]\), i.e.~\(s_2\ge\tfrac52\),
\(s_1+s_4-s_5\ge\tfrac72\), \(s_3\ge\tfrac32\); every
\(s(t)=(\tfrac72-t,\allowbreak{}\tfrac52,\allowbreak{}\tfrac32,\allowbreak{}t,\allowbreak{}0,\allowbreak{}0)\), \(0\le t\le\tfrac12\),
makes all three tight with cost \(\tfrac{15}2\). The endpoints
\(t=0,\allowbreak{}\tfrac12\) (and \(t=\tfrac14\)) satisfy all \(522\) Horn
inequalities of \(d=6\) exactly ({\small\texttt{rank\_\allowbreak{}test.\allowbreak{}py}},
{\small\texttt{nonunique\_\allowbreak{}check.\allowbreak{}py}}, rational arithmetic, two independent Horn
generators), hence so does every \(t\in[0,\allowbreak{}\tfrac12]\) by convexity of
the Horn polyhedron; the LP maximum of \(s_4\) at cost \(\tfrac{15}2\)
is exactly \(\tfrac12\). The layering
\(\mathcal L'_3=\{-\tfrac52\},\allowbreak{}\{-3,\allowbreak{}2\},\allowbreak{}\{2,\allowbreak{}0\},\allowbreak{}\{\tfrac32\}\) with
\(\sigma_2(S_2)=\tfrac14\) (the midpoint of the admissible interval
\([0,\allowbreak{}\tfrac12]\)) can be solved in closed form
({\small\texttt{exact\_\allowbreak{}rank4\_\allowbreak{}matrix.\allowbreak{}py}}):
\(R_1=\begin{pmatrix}p&q\\ q&\frac52-p\end{pmatrix}\) with
\(q^2=p(\frac52-p)\) is rank one, and \(\det(R_1-D_1)=\frac{13}{16}\)
forces \(p=\frac{11}{80}\);
\(R_2=\begin{pmatrix}u&w\\ w&\frac72-u\end{pmatrix}\) with spectrum
\(\{\frac{13}4,\allowbreak{}\frac14\}\) and \(\det(R_2-D_2)=0\) forces
\(u=\frac{99}{32}\), \(w^2=\frac{455}{1024}\). In the basis ordered
\((-b_2,\allowbreak{}-b_1,\allowbreak{}a_1,\allowbreak{}a_2,\allowbreak{}0,\allowbreak{}a_3)\) the resulting real matrix is
\[C=\begin{pmatrix}0&0&0&0&0&0\\[2pt] \frac{\sqrt{110}}{20}&0&0&0&0&0\\[2pt] \frac{3\sqrt{210}}{20}&0&0&0&0&0\\[2pt] 0&\frac{15+91\sqrt{165}}{480}&\frac{39\sqrt{35}-5\sqrt{231}}{480}&0&0&0\\[2pt] 0&\frac{5\sqrt{3003}-3\sqrt{455}}{480}&\frac{15\sqrt{13}+7\sqrt{2145}}{480}&0&0&0\\[2pt] 0&0&0&-\frac{\sqrt{35}}{4}&-\frac{\sqrt{13}}{4}&0\end{pmatrix},\]
numerically \(C_{21}=0.5244\), \(C_{31}=2.1737\), \(C_{42}=2.4665\),
\(C_{43}=0.3224\), \(C_{52}=0.4375\), \(C_{53}=0.7881\),
\(C_{64}=-1.4790\), \(C_{65}=-0.9014\). For this \(C\),
\(CC^{\mathsf T}-C^{\mathsf T}C-2F=0\) \emph{exactly} (verified
symbolically), \(\tfrac12\|C\|^2=\tfrac{15}2=\Phi\),
\(\operatorname{rank}C=4=m+1\), and the common spectrum is
\((\tfrac{13}4,\allowbreak{}\tfrac52,\allowbreak{}\tfrac32,\allowbreak{}\tfrac14,\allowbreak{}0,\allowbreak{}0)\). The exact entries
evaluated in double precision give residual \(1.4\cdot10^{-15}\), and
the floating-point construction of {\small\texttt{my\_\allowbreak{}construct.\allowbreak{}py}} gives
\(4.8\cdot10^{-15}\); the four-decimal display alone gives
\(2.1\cdot10^{-4}\), so the full-precision entries
({\small\texttt{repro/\allowbreak{}data/\allowbreak{}rank4\_\allowbreak{}optimizer\_\allowbreak{}exact.\allowbreak{}json}} and
{\small\texttt{rank4\_\allowbreak{}optimizer\_\allowbreak{}30digits.\allowbreak{}txt}}, 30 significant digits) are
needed to reproduce the residual. The same construction gives
rank-\((m+1)\) optimizers for \(m=5,\allowbreak{}7\) with \(z=1\)
({\small\texttt{rank\_\allowbreak{}more.\allowbreak{}py}}, residuals \(3\cdot10^{-15}\),
\(7\cdot10^{-15}\); e.g.~\(m=5\): \(s=(2.7,\allowbreak{}2.4,\allowbreak{}1.7,\allowbreak{}1.2,\allowbreak{}0.9,\allowbreak{}0.2,\allowbreak{}0,\allowbreak{}0)\)).
For \(z=0\) the same \(\lambda\) without the zero has
\((\mathrm{Pi})_2\) reading \(s_1\ge A_2\) and every optimizer has rank
\(3\).

\textbf{(4) Uniqueness.} By Horn sufficiency the optimal common spectra
are the points of the optimal face of (2.1). \emph{\(m=2\):} in the
\(F_1\) chamber (\(b_2\le a_2\)) the certificate is
\([s_1\ge b_1]+[s_2\ge a_2]\) and tightness gives
\(s=(b_1,\allowbreak{}a_2,\allowbreak{}0,\allowbreak{}\dots)\); in the \(G_1\) chamber (\(b_2\ge a_2\)) it is
\([s_2\ge b_2]+[s_1\ge a_1]\), giving \(s=(a_1,\allowbreak{}b_2,\allowbreak{}0,\allowbreak{}\dots)\); unique
for all \(z\). \emph{\(m=3\):} \(F_2\) (\(b_2\le a_3\)):
\(s=(b_1,\allowbreak{}A_2-b_2,\allowbreak{}a_3,\allowbreak{}0,\allowbreak{}\dots)\). \(F_1\) (\(a_3\le b_2\le a_2\)):
\(s_1=b_1\), \(s_2+s_3=A_2=a_2+a_3\) with \(s_4=\dots=0\); the Weyl
inequalities \(s_2\ge a_2\), \(s_3\ge a_3\) pin
\(s=(b_1,\allowbreak{}a_2,\allowbreak{}a_3,\allowbreak{}0,\allowbreak{}\dots)\). \(F_0\) (\(a_2\le b_2\le E_1-g^*\)):
\(s_1+s_2=P\), \(s_3=a_3\); with \(s_1\ge b_1\) and \(s_2\ge b_2\) this
pins \(s=(b_1,\allowbreak{}b_2,\allowbreak{}a_3,\allowbreak{}0,\allowbreak{}\dots)\). \(G_1\) (\(b_2\ge a_2+a_3\)):
\(s_2=b_2\), \(s_1+s_3=a_1+a_3\), pinned by Weyl to
\((a_1,\allowbreak{}b_2,\allowbreak{}a_3,\allowbreak{}0,\allowbreak{}\dots)\). \(G_3\) with \(z=0\):
\(s=(A_2,\allowbreak{}b_2,\allowbreak{}a_3,\allowbreak{}0,\allowbreak{}0)\). All these pinning inequalities hold in every
\(d\), so the optimal spectrum is unique for \(m=3\) whenever \(z=0\) or
\(b_2\le E_1-a_3\), and (3) shows non-uniqueness in the open \(G_3\)
chamber for \(z\ge1\) (the family \(s(t)\) above, or in general
\((A_2-\sigma,\allowbreak{}b_2,\allowbreak{}a_3,\allowbreak{}\sigma,\allowbreak{}0,\allowbreak{}\dots)\) with
\(\sigma\in[0,\allowbreak{}\min(h,\allowbreak{}a_3)]\)). \emph{\(m\ge4\):} take \(k=m-3\ge1\) and
a point of the open \(F_{m-3}\) chamber
\(a_{m-1}=E_{m-1}<b_2<E_{m-2}=a_{m-2}+a_m\) with \(a_{m-2}>a_{m-1}\). In
\(\mathcal L(F_{m-3})\) the layer after the mixed layer is the pair
\(\{a_{m-2},\allowbreak{}a_{m-1}\}\), followed by the singleton \(\{a_m\}\). Its
state is \((A_{m-2},\allowbreak{}[0,\allowbreak{}h])\) with \(h=\min(b_2,\allowbreak{}A_{m-2}-b_2)>a_{m-1}\),
and the singleton requires \(\tau_2=0\), attainable from
\(\sigma_2=\sigma\) iff
\(\sigma\in[a_{m-1},\allowbreak{}\min(a_{m-2},\allowbreak{}a_{m-1}+a_m)]\). Every
\(\sigma\in[a_{m-1},\allowbreak{}\min(a_{m-2},\allowbreak{}a_{m-1}+a_m,\allowbreak{}h)]\), a nondegenerate
interval, is attainable and gives an optimizer with common spectrum
containing the pair \((A_{m-2}-\sigma,\allowbreak{}\sigma)\); distinct \(\sigma\)
give distinct spectra. (The corresponding \(F_{m-3}\) certificate ends
with \(\Pi_{m-2}\) and \([s_m\ge a_m]\), whose tightness
\(s_{m-2}+s_{m-1}=A_{m-2}\) leaves exactly this one degree of freedom;
the ``free split'' of a certificate alone does not prove non-uniqueness,
the construction does.) For \(m=4\), \(a=(4,\allowbreak{}3,\allowbreak{}2,\allowbreak{}1)\), \(b_2=\tfrac52\):
\(s=(\tfrac{15}2,\allowbreak{}6-\sigma,\allowbreak{}\sigma,\allowbreak{}1,\allowbreak{}0,\allowbreak{}0)\) is Horn-feasible for
\(\sigma=2,\allowbreak{}\tfrac94,\allowbreak{}\tfrac52\) (all \(522\) inequalities,
{\small\texttt{nonunique\_\allowbreak{}check.\allowbreak{}py}}), cost \(\tfrac{29}2=F_1\). LP scans find
non-unique optimal spectra at \(37/60\) random points for \(m=4\),
\(z=0\) (active forms \(F_1,\allowbreak{}G_1,\allowbreak{}G_3\); reviewer) and at \(25/40\)
(\(m=4\)) and \(37/40\) (\(m=6\)) random points (author), at \(3/80\)
and \(2/40\) points for \(m=3\) with \(z=1,\allowbreak{}2\) (all in the \(G_3\)
chamber), and never for \(m=2\) or for \(m=3\), \(z=0\). \(\square\)

\emph{Remark (quantifier).} ``Every optimizer'' refers to HS-optimizers
of \(\tfrac12\|C\|^2\) under \(CC^*-C^*C=2F\). In the two-operator
formulation \(\inf\|A\|\|B\|\), unbalanced product optimizers
\((sA,\allowbreak{}s^{-1}B)\) give \(C=sA+is^{-1}B\) of larger rank already for
\(F=\operatorname{diag}(1,\allowbreak{}-1)\), as noted in the audit of
{[}OneSpike{]}; the rank statements here are for the self-commutator
formulation.

\subsection{8. Consistency checks}\label{consistency-checks}

\textbf{\(m=2\) versus the four-level formula.} On
\(\lambda=(a_1,\allowbreak{}a_2,\allowbreak{}-b_2,\allowbreak{}-b_1)\): \(\lambda_1-\lambda_3=a_1+b_2=G_1\),
\(\lambda_2-\lambda_4=a_2+b_1=F_1\),
\(\lambda_1-2\lambda_2-\lambda_3=G_1-2a_2\),
\(\lambda_2+2\lambda_3-\lambda_4=F_1-2b_2\), and
\(F_0=P\le\max(F_1,\allowbreak{}G_1)\) (Section 6). So \(\Phi=\max(F_1,\allowbreak{}G_1)\)
coincides identically with {[}FourLevel{]} Theorem 3.1; here
\(\kappa_4=\max(\lambda_1-\lambda_3,\allowbreak{}\lambda_2-\lambda_4)\).

\textbf{\(m=3\) versus the twelve five-level terms.} The gap forms
\(q_1,\allowbreak{}\dots,\allowbreak{}q_6\) and their reversals of {[}FiveLevel{]} Theorem 3.1,
written as spectral forms on \(\lambda=(a_1,\allowbreak{}a_2,\allowbreak{}a_3,\allowbreak{}-b_2,\allowbreak{}-b_1)\), are
(sympy, {\small\texttt{symbolic.\allowbreak{}py}}, {\small\texttt{d5\_\allowbreak{}dominated\_\allowbreak{}terms.\allowbreak{}py}}): five of
them are our forms, \(\rho q_5=\lambda_3-\lambda_4-\lambda_5=F_0\),
\(\rho q_2=\lambda_2+\lambda_3-\lambda_5=F_1\),
\(\rho q_4=-\lambda_1+\lambda_3-2\lambda_5=F_2\),
\(q_6=\lambda_1+\lambda_3-\lambda_4=G_1\),
\(\rho q_3=\lambda_2+2\lambda_3-\lambda_4=G_3\); the other seven are
dominated termwise by a single form on the stratum, hence exactly:
\[\begin{aligned}
q_1&=2\lambda_1-2\lambda_2-\lambda_3+\lambda_5=G_1-3a_2-3a_3, & \rho q_1&=-\lambda_1+\lambda_3+2\lambda_4-2\lambda_5=F_2-2b_2,\\
q_2&=\lambda_1-\lambda_3-\lambda_4=G_1-2a_3, & q_3&=\lambda_2-2\lambda_3-\lambda_4=G_3-4a_3,\\
q_4&=2\lambda_1-\lambda_3+\lambda_5=G_1-a_2-3a_3, & q_5&=\lambda_1+\lambda_2-\lambda_3=F_0-3a_3,\\
\rho q_6&=\lambda_2-\lambda_3-\lambda_5=F_1-2a_3.
\end{aligned}\] ({\small\texttt{verify\_\allowbreak{}mn2.\allowbreak{}py}} PART3 also checks both formulas
against \(\Phi\) at \(3000\) exact stratum points each.)

\textbf{One-spike limit.} At \(b_2=0\), \(E_{m+1}=0\le b_2\le E_m\)
selects \(F_{m-1}=(m-1)P+\sum_{j<m}(j-m+1)a_j+a_m=\sum_j j\,\allowbreak{}a_j\), the
one-spike formula of {[}OneSpike{]} Theorem 3.1 after \(F\mapsto-F\);
the certificate and the layering \(\mathcal L(F_{m-1})\) remain valid at
\(b_2=0\), and the extension is also forced by continuity of
\(\kappa_d\) on the chamber ({[}FourLevel{]} Theorem 2.2, {[}Ceiling{]}
Section 2).

\textbf{Reduction \(a_m\to0\).} At \(a_m=0\) the forms of \(\Phi_m\)
reduce to those of \(\Phi_{m-1}\) except \(F_{m-1}\), which becomes
\(F_{m-2}-b_2<F_{m-2}\), and (for \(m\) odd) \(G_m\), which becomes
\(G_{m-2}-g_{m-2}\le G_{m-2}\); so \(\Phi_m|_{a_m=0}=\Phi_{m-1}\), as
padding invariance requires.

\textbf{Padding.} The example of {[}RankOnset{]} Remark 3.3,
\((25,\allowbreak{}18,\allowbreak{}18,\allowbreak{}-10,\allowbreak{}-17,\allowbreak{}-17,\allowbreak{}-17)/7\) with \(\kappa_7=74/7\ne\kappa_8=73/7\),
has inertia \((3,\allowbreak{}4)\) and lies outside both the \((m,\allowbreak{}2)\) stratum and
its mirror \((2,\allowbreak{}n)\); it does not contradict Theorem A. On our stratum
the value is padding-invariant, while Section 7 shows that the optimizer
set is not.

\subsection{9. Exact verification}\label{exact-verification}

All checks use rational arithmetic unless stated otherwise.

\emph{Author's scripts} ({\small\texttt{verify\_\allowbreak{}mn2.\allowbreak{}py\ 8}}, {\small\texttt{lr.\allowbreak{}py}},
{\small\texttt{layer\_\allowbreak{}chain.\allowbreak{}py\ 13\ 300}}). {\small\texttt{lr.\allowbreak{}py}}: the criterion
\(c^{\lambda(K)}_{\lambda(I)\lambda(J)}>0\) reproduces Fulton's
recursive \(T^d_r\) exactly for \(d\le7\) (counts
\(3;\allowbreak{}\ 6,\allowbreak{}6;\allowbreak{}\ 10,\allowbreak{}21,\allowbreak{}10;\allowbreak{}\ 15,\allowbreak{}56,\allowbreak{}56,\allowbreak{}15;\allowbreak{}\ 21,\allowbreak{}126,\allowbreak{}228,\allowbreak{}126,\allowbreak{}21;\allowbreak{}\ 28,\allowbreak{}252,\allowbreak{}751,\allowbreak{}751,\allowbreak{}252,\allowbreak{}28\))
and gives \(c=1\) for the Pieri triples for \(d\le12\). PART1: for every
form, \(m\le8\), \(z\in\{0,\allowbreak{}1,\allowbreak{}2\}\) (\(675\) triples), the certificate
list is generated and checked for the sum condition, \(c\ge1\),
membership in \(T^d_r\) for \(d\le8\), all left coefficients \(\le1\),
and right side \(=\) form. PART2: at \(400\) exact points per \(m\le8\),
feasibility of each layering \(\iff\) its closed-form region, the argmax
layering is feasible, chambers as in Theorem B. PART2b: explicit \(C\)
at \(60\) points per \(m\le8\) with residuals \(\le2\cdot10^{-8}\)
(float eigensolves at ties; \(10^{-14}\) elsewhere) and rank \(m\).
PART3: the \(d=4\) and \(d=5\) formulas at \(3000\) points each; the
one-spike limit for \(m\le8\). {\small\texttt{layer\_\allowbreak{}chain.\allowbreak{}py}}: \(300\) random
exact points per \(m\le13\), no point without a feasible layering,
histogram of the number of feasible layerings.

\emph{Reviewer's scripts} (independent implementations; the author's
code was imported only for comparison). {\small\texttt{my\_\allowbreak{}horn.\allowbreak{}py}}: own Fulton
generator and own LR coefficients (bialternant formula, Jacobi--Trudi),
agreeing on the counts above. {\small\texttt{my\_\allowbreak{}cert.\allowbreak{}py}}: all \(168\)
certificates for \(m\le7\), \(z\in\{0,\allowbreak{}1,\allowbreak{}2,\allowbreak{}5\}\), with a different
elimination of \(b_1\); the systematic scan lists exactly the nine
certificates with a \(+s_{m+1}\) term (\(G_3\), \(G_5\), \(G_7\) at
\(z=1,\allowbreak{}2,\allowbreak{}5\)). {\small\texttt{my\_\allowbreak{}construct.\allowbreak{}py}}, {\small\texttt{adversarial.\allowbreak{}py}}: own
interval recursion agreeing with the author's on \(2708\) exact
boundary-heavy points (\(b_1=b_2\), every \(b_2=E_k\), every
\(b_2=E_1-g_j\), flat \(a\), ties, \(a_m=10^{-6}\), one dominant
\(a_1\), geometric \(a\), \(m\le12\)), no point without a feasible
layering; \(732\) explicit matrices at boundary points with residuals
\(\le8.4\cdot10^{-9}\) and rank \(m\). {\small\texttt{symbolic.\allowbreak{}py}}: all
identities of Theorem B and of Section 8 symbolically for \(m\le9\).
{\small\texttt{lp\_\allowbreak{}checks.\allowbreak{}py}}: an own float Horn LP at \(600\) random points
for each of nine \((m,\allowbreak{}z)\) pairs (\(5400\) points, \(d\le7\)) agrees
with \(\Phi\) to \(10^{-9}\) and returns exactly the forms \(F_k,\allowbreak{}G_j\)
as dual gradients; uniqueness scans as quoted in Section 7.
{\small\texttt{rank\_\allowbreak{}test.\allowbreak{}py}}, {\small\texttt{rank\_\allowbreak{}more.\allowbreak{}py}}: the exact rank-\(4\)
certificate and the rank-\((m+1)\) optimizers for \(m=3,\allowbreak{}5,\allowbreak{}7\);
exposed-form counts. The hive-LP duals of {\small\texttt{mn2\_\allowbreak{}forms.\allowbreak{}py}} (m
\(\le9\), \(z\le2\)) were the original source of the conjectured forms.

\emph{Scripts added in version 2} (all in {\small\texttt{repro/\allowbreak{}}}, run by
{\small\texttt{run\_\allowbreak{}all.\allowbreak{}py}}). {\small\texttt{exact\_\allowbreak{}rank4\_\allowbreak{}matrix.\allowbreak{}py}}: the rank-\(4\)
optimizer of Section 7 in closed form (sympy), residual exactly zero,
entries written to {\small\texttt{data/\allowbreak{}}} with 30 significant digits.
{\small\texttt{coverage\_\allowbreak{}exact.\allowbreak{}py}}: for \(2\le m\le10\), on a rational grid of
\(9567\) points containing every breakpoint \(E_k\), \(E_1-g_j\) and
\(b_1=b_2\), flat, tied, geometric and dominant-gap spectra, exact
checks of cost \(=\) form, feasibility \(\iff\) region for every
layering (Proposition 5.4), coverage (Proposition 5.5), the chambers and
comparison identities of Theorem B, the feasibility and rank-\((m+1)\)
criteria of Theorem C\('\)(3) for odd \(m\), the nondegenerate family of
Theorem C\('\)(4) for \(m\ge4\), and the corrected activity criterion
for the \(G\)-forms of Section 6. {\small\texttt{pieri\_\allowbreak{}triples\_\allowbreak{}search.\allowbreak{}py}}:
exhaustive enumeration of \(T^d_r\) (\(d\le9\)) of all Horn triples
sharing the \(s\)-form of \((\mathrm{Pi})_i\), and the
Lidskii--Wielandt-only LP comparison of Section 4.1. The workshop's own
script {\small\texttt{check\_\allowbreak{}horn.\allowbreak{}py}} (Fulton recursion up to \(d=7\), \(200\)
float LPs, exact Horn checks of the six spectra of Section 7) was rerun
on the author's machine and reproduces its published output (maximal LP
discrepancy \(7.1\cdot10^{-15}\); log in
{\small\texttt{logs/\allowbreak{}taller\_\allowbreak{}check\_\allowbreak{}horn.\allowbreak{}json}}).

The received version reports that its reruns
({\small\texttt{repro/\allowbreak{}run\_\allowbreak{}all.\allowbreak{}py}}) reproduce the outputs quoted above; its
reported runtimes are listed in {\small\texttt{repro/\allowbreak{}README.\allowbreak{}md}}. The present
workshop's own reruns and their limits are itemized separately in the
v003 revision report and verification records; the full runner was not
repeated in this revision.

\subsection{10. Open questions}\label{open-questions}

\begin{enumerate}
\def\labelenumi{\arabic{enumi}.}
\tightlist
\item
  \emph{Inertia \((m,\allowbreak{}3)\) and beyond.} The certificate here uses one
  Pieri-type triple per even index; a natural guess is that inertia
  \((m,\allowbreak{}n)\) requires triples adding \(n-1\) boxes to a rectangle, and
  that block shifts with blocks of size up to \(n\) attain. Neither is
  known; the \(d=8\) example of {[}RankOnset{]} with \(r_*>r_0\) has
  inertia \((4,\allowbreak{}4)\) and shows that the rank-\(m\) block picture cannot
  persist unchanged.
\item
  \emph{A general interior formula} for \(\kappa_d\) on the ordered
  chamber remains open beyond \(d=5\); the present result is the first
  infinite family of strata with a closed formula in which both signs
  have multiplicity \(>1\).
\item
  \emph{Balanced stability constants \(\gamma_N\).} The optimal forward
  coefficients of {[}Ceiling,Bal44,Bal55{]} are known exactly for
  \(N=2,\allowbreak{}\dots,\allowbreak{}5\) (\(1/2\), \(17/36\), \(4/7\), \(113/152\)). Hive-LP
  computations reported in the author's line notes give the
  \emph{numerical} values \(\gamma_6=9/10\), \(\gamma_7=781/740\),
  \(\gamma_8=317/264\), \(\gamma_9=23/17\), \(\gamma_{10}=3/2\), with
  the extremal family switching at \(N=7\) from the \((A_N,\allowbreak{}B)\) pairs to
  a \((u_3,\allowbreak{}z_N)\) family; these are numerical only, without exact
  certificates.
\item
  \emph{Classification of optimizers} for \(m\ge4\): which chambers have
  a unique optimal spectrum, and a description of the optimal face.
\item
  \emph{Which Horn triples are needed in general?} On this stratum all
  facets are certified by Lidskii--Wielandt, Weyl and a single Pieri
  family; whether the facets of \(\kappa_d\) on every stratum admit
  certificates with LR coefficient \(1\) and bounded ``distance from
  Lidskii'' is open.
\end{enumerate}

\subsection{11. Reproducibility}\label{reproducibility}

The directory {\small\texttt{repro/\allowbreak{}}} accompanying this manuscript contains the
author's scripts ({\small\texttt{lr.\allowbreak{}py}}, {\small\texttt{layer\_\allowbreak{}chain.\allowbreak{}py}},
{\small\texttt{verify\_\allowbreak{}mn2.\allowbreak{}py}}, {\small\texttt{hive\_\allowbreak{}core.\allowbreak{}py}},
{\small\texttt{check\_\allowbreak{}horn\_\allowbreak{}lp.\allowbreak{}py}}, {\small\texttt{mn2\_\allowbreak{}forms.\allowbreak{}py}},
{\small\texttt{conj\_\allowbreak{}formula.\allowbreak{}py}}, {\small\texttt{horn\_\allowbreak{}duals\_\allowbreak{}pad.\allowbreak{}py}}), the reviewer's
scripts ({\small\texttt{my\_\allowbreak{}horn.\allowbreak{}py}}, {\small\texttt{my\_\allowbreak{}cert.\allowbreak{}py}},
{\small\texttt{my\_\allowbreak{}construct.\allowbreak{}py}}, {\small\texttt{rank\_\allowbreak{}test.\allowbreak{}py}},
{\small\texttt{rank\_\allowbreak{}more.\allowbreak{}py}}, {\small\texttt{lp\_\allowbreak{}checks.\allowbreak{}py}}, {\small\texttt{adversarial.\allowbreak{}py}},
{\small\texttt{symbolic.\allowbreak{}py}}) with their logs, two small scripts written for
this manuscript ({\small\texttt{nonunique\_\allowbreak{}check.\allowbreak{}py}},
{\small\texttt{d5\_\allowbreak{}dominated\_\allowbreak{}terms.\allowbreak{}py}}), the three scripts added in version 2
({\small\texttt{exact\_\allowbreak{}rank4\_\allowbreak{}matrix.\allowbreak{}py}}, {\small\texttt{coverage\_\allowbreak{}exact.\allowbreak{}py}},
{\small\texttt{pieri\_\allowbreak{}triples\_\allowbreak{}search.\allowbreak{}py}}) with the data files they write in
{\small\texttt{data/\allowbreak{}}}, the workshop script {\small\texttt{check\_\allowbreak{}horn.\allowbreak{}py}} and its
output, a runner {\small\texttt{run\_\allowbreak{}all.\allowbreak{}py}}, and {\small\texttt{README.\allowbreak{}md}} with
commands, measured runtimes and the expected final lines of every
script. Dependencies: Python 3, {\small\texttt{numpy}}, {\small\texttt{scipy}} (HiGHS),
{\small\texttt{sympy}}. Literature searches (2026-09-07, two queries on the
inverse self-commutator problem with prescribed spectrum and on block
weighted shifts / Horn certificates) returned only generic material on
almost-commuting matrices, commutator norm inequalities and Horn
expositions.

\subsection{AI-assistance statement}\label{ai-assistance-statement}

The received version declares substantial generative-AI assistance in
drafting, mathematical derivations, code and review, attributed there to
``Claude Fable 5.1'' (Anthropic). That historical model identity, the
reported review sessions and their independence have not been
authenticated by this workshop. Historical descriptions of an ``author''
or ``reviewer'' script identify the supplied files, not a verified
independent editorial review. OpenAI Codex assisted with the present
private workshop revision: checking arguments and sources, correcting
identified errors, running the checks recorded in the accompanying
revision report, and preparing the PDF. Numerical checks support only
their tested instances. This assistance is not an independent human peer
review or a formal AIRR model-assessment report. Lluis Eriksson is the
declared author and AIRR founder/operator; this conflict must be handled
by the separate editorial process. Scientific authorship and approval
remain with the author.

The deposited v003 was subsequently assessed in actual OpenAI Codex
calls by gpt-6-astra and gpt-5.6-sol on 9 September 2026. Their original
reports and subsequent clarifications are retained. Both sets of
comments informed this v004 revision; Astra had also participated in
earlier manuscript preparation. A separate reviewing context does not
erase that history. Neither the v003 reports nor this disclosure
constitute approval of the revised v004 artifact.

\subsection{Changes in internal revision v004 (9 September
2026)}\label{changes-in-internal-revision-v004-9-september-2026}

\begin{enumerate}
\def\labelenumi{\arabic{enumi}.}
\tightlist
\item
  Expanded the third candidate inequality in Proposition 5.4(3) into its
  exact sum of ordered differences, including empty sums and equality
  cases. The general identity has a direct derivation and supplementary
  symbolic checks.
\end{enumerate}

Expanded the induction for both interval endpoints in Proposition 5.4(3)
and supplied the complete d=6, i=2 Horn/Pieri example, including
partitions, coefficient and both sides of the inequality.

The earlier deposited v003 and its model reports are retained. This
revision responds to their findings; it is not a final acceptance.

\subsection{Changes in internal revision v003 (8 September
2026)}\label{changes-in-internal-revision-v003-8-september-2026}

\begin{enumerate}
\def\labelenumi{\arabic{enumi}.}
\tightlist
\item
  Corrected the exact criterion for an empty open \(F_k\) chamber:
  equality within each indicated pair, not constancy of the entire tail.
\item
  Added the exact counterexample \((4,\allowbreak{}3,\allowbreak{}3,\allowbreak{}1,\allowbreak{}1)\) and the nonnegative-sum
  proof of the corrected criterion.
\item
  Made the unpadded case an explicit alternative in Theorem C'(2),
  rather than implying that it forces a different active form.
\item
  Clarified the rank argument and qualified historical AI-review
  provenance. The main cost formula is unchanged.
\item
  Distinguished the operator-norm objective in {[}JOS13{]} from the
  mixed operator/Hilbert--Schmidt objective in {[}AS15{]}, and completed
  the latter's bibliographic entry.
\end{enumerate}

\subsection{Changes in version 2 (8 September
2026)}\label{changes-in-version-2-8-september-2026}

Version 2 answers the items of the external workshop evaluation
(TALLER-0004); the workshop's page numbers refer to the version-1 PDF.

\begin{enumerate}
\def\labelenumi{\arabic{enumi}.}
\tightlist
\item
  \emph{(Item: page 9, Section 6, identity
  \(E_{k+1}-E_{k+2}=a_{k+1}\).)} Accepted: the identity was false
  (\(a=(4,\allowbreak{}3,\allowbreak{}2,\allowbreak{}1)\), \(k=1\) gives \(E_2-E_3=2\ne3\)). The exposedness
  argument now uses the correct identities \(E_{k+1}-E_{k+3}=a_{k+1}\)
  and \(E_{k+1}-E_{k+2}=\sum(\text{pair gaps})\,\allowbreak{}(+a_m)\), characterises
  exactly when the open \(F_k\) chamber is empty, and exhibits explicit
  exposing spectra. Theorems A, B, C\('\) are unchanged.
\item
  \emph{(Item: page 6, Weyl row \((3;\allowbreak{}3;\allowbreak{}3)\).)} Accepted: the triple is
  \((3;\allowbreak{}1;\allowbreak{}3)\) (\(I=K=\{3\}\), \(J=\{1\}\)), as in Lemma 4.1 and in the
  reviewer's {\small\texttt{my\_\allowbreak{}cert.\allowbreak{}log}}; the Weyl item of Lemma 4.1 now
  states \(I,\allowbreak{}J,\allowbreak{}K\) explicitly. No other certificate display was
  affected.
\item
  \emph{(Item: last sentence of Section 6.)} Accepted: the sentence is
  replaced by the exact criterion ``some \(G_j\) is active somewhere on
  the stratum iff \(\sum_{j\ne j^*}g_j\le g^*\), uniquely iff strict'',
  with the tie cases (including the workshop's example \(a=(2,\allowbreak{}1,\allowbreak{}1)\),
  \(b_1=b_2=2\), \(F_0=G_1=G_3=5\)) spelled out.
\item
  \emph{(Item: page 11 matrix at full precision.)} Accepted: the
  rank-\(4\) optimizer is now given in closed form (square roots of
  integers), with residual exactly zero; the workshop's residual
  \(2.15\cdot10^{-4}\) from the four-decimal print is reproduced and
  explained; data files with 30 significant digits and the script are in
  {\small\texttt{repro/\allowbreak{}}}.
\item
  \emph{(Item: re-audit of Sections 5.2--5.3 and of Theorem C\('\).)} No
  mathematical error was found in Lemmas 5.1--5.3, Proposition 5.4 or
  Proposition 5.5, but the exposition was tightened: a new Lemma 5.0
  states why feasibility reduces exactly to the interval recursion; the
  proofs of Proposition 5.4 list every tail case (\(k=m-1\); \(j=1\),
  \(j=m-1\), \(j=m\); \(m\) even/odd; \(m=2\)); the rank bookkeeping in
  the proof of Theorem C\('\)(3) was corrected (version 1 wrote
  \(\operatorname{rank}C=1+2\#\{\dots\}+1\); the correct count is
  \(\operatorname{rank}C=m+[\sigma_2(S_{T-1})>0]\), because every
  earlier \(S_t\) has rank \(2\) in any solution), a misleading
  parenthetical in the proof of (1)--(2) was removed, and the
  exhaustiveness of cases (2)/(3) is stated. A new exact-arithmetic
  script ({\small\texttt{coverage\_\allowbreak{}exact.\allowbreak{}py}}) checks all these statements on a
  fine rational grid with all ties and boundaries for \(m\le10\).
\item
  \emph{(Item: dependence on internal record codes.)} The antecedent
  records are cited by descriptive keys with theorem numbers and are
  restated where used; the reference list gives their titles and ARR
  identifiers. The sentence citing a script ({\small\texttt{find\_\allowbreak{}triples.\allowbreak{}py}})
  that was not part of the reproducibility package was replaced by a
  verifiable statement backed by {\small\texttt{pieri\_\allowbreak{}triples\_\allowbreak{}search.\allowbreak{}py}}.
\item
  The workshop's {\small\texttt{check\_\allowbreak{}horn.\allowbreak{}py}} was rerun and agrees with its
  published results (Section 9). The reproducibility package was
  completed (README with expected outputs, data files, logs).
\end{enumerate}

\subsection{References}\label{references}

\begin{itemize}
\tightlist
\item
  {[}FourLevel{]} L. Eriksson, \emph{The Exact Four-Level Inverse
  Commutator Cost: Horn--Littlewood--Richardson Facets, Rank
  Transitions, and Sharp Loop Synthesis}, AI Research Record,
  ARR-2026-3M1EEG1T689ADSMW (2026). Cited: Lemma 2.1 (two-operator
  reduction), Theorem 2.2 (Horn program), Theorem 3.1 (four-level
  formula).
\item
  {[}FiveLevel{]} L. Eriksson, \emph{The Exact Five-Level Inverse
  Commutator Cost: Twelve Horn Chambers, Optimal Rank, and the Sharp 5/2
  Resource Tax}, AI Research Record, ARR-2026-37B8R0QTA894GTFF (2026).
  Cited: Theorem 3.1 (five-level formula).
\item
  {[}RankAdaptive{]} L. Eriksson, \emph{Sharp Rank-Adaptive Bounds for
  Inverse Self-Commutators}, AI Research Record,
  ARR-2026-1D2QV1RP1292JREW (2026). Cited: Theorem 1.1 (rank-adaptive
  bound), the weighted-shift construction of Section 4.
\item
  {[}OneSpike{]} L. Eriksson, \emph{One-Spike Inverse Self-Commutators
  and Exact Three-versus-Four-Kick Curvature Synthesis}, AI Research
  Record, ARR-2026-7NPRNBW4488HG90K (2026). Cited: Theorem 3.1
  (one-spike formula and rigidity).
\item
  {[}RankOnset{]} L. Eriksson, \emph{Sharp Onset and Unbounded Growth of
  Norm-Optimal Self-Commutator Rank}, AI Research Record,
  ARR-2026-5QQF95VHTC9GABH8 (2026). Cited: Lemma 2.1 (inertia
  obstruction), Proposition 2.2 (Horn program in the
  \(s\)-normalisation), Theorem 3.1 (\(r_*=r_0\) for \(d\le7\)), Theorem
  4.1 (the \(d=8\) phase), Remark 3.3 (padding example).
\item
  {[}Ceiling{]} L. Eriksson, \emph{Sharp inertia ceilings and optimal
  stability for inverse self-commutators}, AI Research Record,
  ARR-2026-24M24KDPZK8HDBQ9 (2026). Cited: Theorem 1 (inertia ceiling),
  Section 2 (continuity), Theorem 3 (balanced stability constants).
\item
  {[}Bal44{]} L. Eriksson, \emph{Polynomial vertex reduction and optimal
  stability at balanced inertia (4,4)}, AI Research Record,
  ARR-2026-54Q3HMFJ0Z8CZB4T (2026).
\item
  {[}Bal55{]} L. Eriksson, \emph{Optimal stability and all equality
  cases at balanced inertia (5,5)}, AI Research Record,
  ARR-2026-6ZJY1SSSJA98WBWX (2026).
\item
  {[}Ho62{]} A. Horn, Eigenvalues of sums of Hermitian matrices,
  \emph{Pacific J. Math.} 12 (1962) 225--241.
\item
  {[}Kl98{]} A. A. Klyachko, Stable bundles, representation theory and
  Hermitian operators, \emph{Selecta Math. (N.S.)} 4 (1998) 419--445.
\item
  {[}KT99{]} A. Knutson, T. Tao, The honeycomb model of
  \(GL_n(\mathbb C)\) tensor products I: proof of the saturation
  conjecture, \emph{J. Amer. Math. Soc.} 12 (1999) 1055--1090.
\item
  {[}Fu00{]} W. Fulton, Eigenvalues, invariant factors, highest weights,
  and Schubert calculus, \emph{Bull. Amer. Math. Soc.} 37 (2000)
  209--249.
\item
  {[}Fu97{]} W. Fulton, \emph{Young Tableaux}, LMS Student Texts 35,
  Cambridge University Press, 1997 (Pieri rule).
\item
  {[}Li50{]} V. B. Lidskii, On the characteristic numbers of the sum and
  product of symmetric matrices, \emph{Dokl. Akad. Nauk SSSR} 75 (1950)
  769--772.
\item
  {[}Wi55{]} H. Wielandt, An extremum property of sums of eigenvalues,
  \emph{Proc. Amer. Math. Soc.} 6 (1955) 106--110.
\item
  {[}We12{]} H. Weyl, Das asymptotische Verteilungsgesetz der Eigenwerte
  linearer partieller Differentialgleichungen, \emph{Math. Ann.} 71
  (1912) 441--479.
\item
  {[}JOS13{]} W. B. Johnson, N. Ozawa, G. Schechtman, A quantitative
  version of the commutator theorem for zero trace matrices, \emph{Proc.
  Natl. Acad. Sci. USA} 110 (2013) 19251--19255.
\item
  {[}AS15{]} O. Angel, G. Schechtman, The Hilbert--Schmidt version of
  the commutator theorem for zero trace matrices, \emph{Bull. London
  Math. Soc.} 47 (2015), no. 4, 715--719; doi:10.1112/blms/bdv045;
  arXiv:1503.07980.
\end{itemize}

\end{document}
